\documentclass[11pt,a4paper]{article}
\usepackage[margin=1in]{geometry}
\usepackage[T1]{fontenc}
\usepackage{lmodern}
\usepackage{microtype}
\usepackage{amsmath,amssymb,amsthm,mathtools}
\usepackage{enumitem}
\usepackage{xcolor}
\usepackage{booktabs,tabularx,array}
\usepackage[numbers,sort&compress]{natbib}
\usepackage{xurl}
\usepackage{hyperref}
\hypersetup{colorlinks=true,linkcolor=blue!45!black,citecolor=blue!45!black,
 urlcolor=blue!45!black,pdftitle={Sample Complexity and Communication Limits of Decentralized Optimization under Relaxed Smoothness},
 pdfauthor={Colombo Team}}
\setlength{\parindent}{1.2em}
\setlength{\parskip}{0.18em}
\setlist[itemize]{leftmargin=*,itemsep=0.25em,topsep=0.35em}
\setlist[enumerate]{leftmargin=*,itemsep=0.25em,topsep=0.35em}
\newcommand{\R}{\mathbb{R}}
\newcommand{\E}{\mathbb{E}}
\newcommand{\eps}{\varepsilon}
\newcommand{\norm}[1]{\left\lVert #1\right\rVert}
\newcommand{\inner}[2]{\left\langle #1,#2\right\rangle}
\newcommand{\cO}{\mathcal{O}}
\newcommand{\cOmega}{\Omega}
\newcommand{\tO}{\widetilde{\mathcal{O}}}
\newcommand{\tOmega}{\widetilde{\Omega}}
\newcommand{\tTheta}{\widetilde{\Theta}}
\newcommand{\one}{\mathbf{1}}
\DeclareMathOperator*{\argmin}{arg\,min}
\DeclareMathOperator*{\argmax}{arg\,max}
\newtheorem{theorem}{Theorem}[section]
\newtheorem{lemma}[theorem]{Lemma}
\newtheorem{proposition}[theorem]{Proposition}
\newtheorem{corollary}[theorem]{Corollary}
\theoremstyle{remark}
\newtheorem{remark}[theorem]{Remark}
\newcolumntype{Y}{>{\raggedright\arraybackslash}X}
\emergencystretch=2em
\allowdisplaybreaks[1]
\title{Sample Complexity and Communication Limits\\
of Decentralized Optimization\\under Relaxed Smoothness}
\author{Colombo Team\\
\small\href{mailto:colombo.2026@outlook.com}{\texttt{colombo.2026@outlook.com}}}
\date{}
\begin{document}
\maketitle
\begin{abstract}
We study decentralized nonconvex optimization when local gradient variation
is controlled by a constant plus the local gradient norm. Each of $m$ agents
has access to an unbiased stochastic gradient oracle with fresh noise of
variance at most $\sigma^2$, and local gradients differ from their average
by at most $H$. Writing $L_H=L_0+L_1H$, we characterize the minimax total
number of local oracle calls, up to logarithmic factors, throughout the
range $\eps^2\le c\Delta L_H$:
\[
\widetilde\Theta\!\left(
 \left[\frac{L_H\Delta}{\eps^2}+\frac{L_1\Delta}{\eps}\right]
       \left[1+\frac{\sigma^2}{\eps^2}\right]
 +\frac{L_H\Delta}{\eps^2}
       \sqrt{1+\min\{m,H^2/\eps^2\}}
 +1+\min\{m,H^2/\eps^2\}\right).
\]
The upper bound uses robust dyadic gradient differences to separate
component uncertainty from fresh oracle noise. The lower bounds
combine an orthogonal embedding, hidden-sign discovery, a retiring
finite-sum chain with uniformly bounded disagreement, and a terminally flat
stochastic chain. The last construction shows that fresh noise multiplies
the coarse-accuracy scale $L_1\Delta/\eps$. The lower bounds allow arbitrary
randomized adaptive queries in sufficiently large finite dimension.
We also distinguish total sample complexity from communication complexity.
For a known fixed gossip matrix $W$, exact arithmetic permits averaging in
$r(W)-1$ communication calls, where $r(W)$ counts its distinct eigenvalues.
Small heterogeneity additionally permits communication-free optimization.
These facts delimit the scope of spectral-gap-based communication bounds.
\end{abstract}

\noindent\textbf{Keywords.} Decentralized optimization; nonconvex optimization;
relaxed smoothness; oracle complexity; variance reduction; communication complexity.
\par\medskip
\begin{samepage}
\noindent\textbf{AI Usage.}
Nearly the entire research pipeline for this paper was carried out by
Colombo, powered by GPT-5.6 Sol.
\par
\end{samepage}

\section{Introduction}

Decentralized optimization separates access to data from the exchange of
information. An agent can sample its own objective without communicating,
but its local gradient may differ substantially from the gradient of the
network average. Under relaxed smoothness, this disagreement has a second
effect: it changes the curvature that an algorithm must accommodate.
A useful complexity theory must therefore account for the interaction
between accuracy, heterogeneity, stochastic noise, and the number of agents.

We consider local functions satisfying the finite-step
$(L_0,L_1)$-smoothness condition, under which gradient variation can grow
with the gradient norm. Luo et al.~\cite{relaxed2025} established a
convergence framework for decentralized normalized stochastic gradient
descent in this setting and identified the inherited smoothness parameter
$L_H=L_0+L_1H$. Their results motivate two distinct questions. Which
combinations of $L_H$, the initial gap $\Delta$, the target accuracy $\eps$,
and the noise level are unavoidable in the \emph{total} number of local
queries? Separately, what does a spectral-gap dependence say when the
communication matrix is fixed and exact linear arithmetic is allowed?

Our answer to the first question is a minimax characterization over all
finite dimensions. The rate has four contributions: ordinary smooth
stochastic optimization, sequential component identification, one-time
component discovery, and stochastic progress at coarse accuracy. These
contributions have different constructions and become dominant in different
parameter ranges. The distinction is particularly useful when local oracle
noise is small, or when the target accuracy is larger than $L_H/L_1$.

The analysis develops three complementary arguments.
\begin{itemize}
\item \textbf{Separating the two sources of uncertainty.}
A dyadic estimator averages component differences and fresh endpoint noise
at different scales. A robust aggregation rule gives simultaneous control
along the adaptive path, including the random times at which the algorithm
returns to a macro-iterate. This yields an additive fresh-noise cost rather
than a product between noise and the effective population size.
\item \textbf{Realizing the lower-bound scales in the stated class.}
An orthogonal simplex makes ordinary $L_H$-smooth hard instances admissible
under local $(L_0,L_1)$-smoothness. A retiring selector adapts dense weak
hiding~\cite{dense2026} so that gradient disagreement does not accumulate
with the number of stages. A terminally flat chain combines low-amplitude
hidden gradients with a pointwise smoothness certificate and a fresh
rare-reveal oracle.
\item \textbf{Specifying the communication quantifiers.}
Classical finite-time consensus~\cite{finiteconsensus2007} supplies a
fixed-matrix implementation of the sample-efficient method. A two-agent
example and a communication-free regime show why an uncapped
$\gamma^{-1/2}$ expression cannot be a universal pointwise communication
lower bound in this exact-real model. These observations are compatible
with lower bounds for hard network families under different protocol
restrictions~\cite{ludesa2021}.
\end{itemize}

Section~\ref{sec:model} specifies the oracle model and states the main
result. Section~\ref{sec:related} compares the result with prior work using
consistent sample counts and explicit assumptions.
Sections~\ref{sec:upper-main} and~\ref{sec:lower-main} give the upper-bound
framework, the main reductions, and the assembly of the minimax rate.
Section~\ref{sec:communication} treats communication separately.
The appendices contain the adaptive-estimation and sequential-hiding
arguments, together with a one-dimensional deterministic lower bound.

\section{Oracle model and main result}
\label{sec:model}

\subsection{Functions, queries, and outputs}

Fix $m\ge2$, $L_0,L_1,\Delta>0$, and $H,\sigma\ge0$.
For a finite dimension $d$, consider continuously differentiable functions
$f_i:\R^d\to\R$ with average
\[
 f(x)=\frac1m\sum_{i=1}^m f_i(x),\qquad
 f(x_0)-\inf_x f(x)\le\Delta.
\]
Only the average objective is required to be bounded below. For every agent
$i$ and every pair with $\norm{x-y}\le1/L_1$, assume
\begin{align}
 \norm{\nabla f_i(x)-\nabla f_i(y)}
 &\le\bigl(L_0+L_1\norm{\nabla f_i(x)}\bigr)\norm{x-y},
 \label{eq:local-relaxed}\\
 \norm{\nabla f_i(x)-\nabla f(x)}&\le H.
 \label{eq:heterogeneity}
\end{align}
The heterogeneity bound is uniform over agents and over the whole space.

A local query specifies an agent $i$ and a point $x$ and returns a gradient
vector $G_i(x;\xi)$. If $\mathcal F$ denotes the complete history before the
call, its conditional moments satisfy
\begin{equation}
 \E[G_i(x;\xi)\mid\mathcal F]=\nabla f_i(x),\qquad
 \E[\norm{G_i(x;\xi)-\nabla f_i(x)}^2\mid\mathcal F]\le\sigma^2.
 \label{eq:oracle-model}
\end{equation}
Every call uses fresh randomness. In particular, the interface does not
allow a random sample to be reused at a second query point, and imposes no
smoothness condition on the noisy response as a function of $x$.
Calls may be serialized, so their noise vectors form martingale differences.
The oracle returns no function values. Algorithms may select both the
agent and the query point adaptively, with unrestricted local computation.

Communication uses a known fixed matrix $W\in\R^{m\times m}$ that is
symmetric, nonnegative, and doubly stochastic, with
\[
 0\preceq W\preceq I,\qquad
 \ker(I-W)=\operatorname{span}\{\one\}.
\]
One communication call applies $W$ to any finite stack of local vectors.
Exact-real local arithmetic is allowed, including linear combinations with
coefficients determined by the known matrix. Write
\begin{equation}
 \gamma=1-\lambda_2(W),\qquad
 r(W)=\#\{\text{distinct eigenvalues of }W\}.
 \label{eq:network-parameters}
\end{equation}
All agents know $x_0$, the parameter bounds, and the protocol, and may use
public randomness. Each agent must return a locally available point
$\widehat x_i$ such that
\begin{equation}
 \E\norm{\nabla f(\widehat x_i)}\le\eps.
 \label{eq:output-criterion}
\end{equation}
The outputs need not be equal; the upper-bound algorithm below produces a
common output. Let $S$ count total local oracle calls across all agents and
let $C_{\rm comm}$ count communication calls.

For each $d$, let $S_d^*$ be the infimum of the worst-case expected value of
$S$ over randomized adaptive protocols satisfying
\eqref{eq:output-criterion} on the entire $d$-dimensional class. Define
\begin{equation}
 S^*:=\sup_{d\ge1}S_d^*.
 \label{eq:minimax-definition}
\end{equation}
Thus the upper bounds are dimension independent, whereas the lower bounds
may choose a sufficiently large finite dimension. The matrix $W$ is held
fixed throughout this definition. Communication is permitted and counted
separately; it is not constrained when minimizing $S$.

\subsection{Scales and notation}

Define
\begin{align}
 L_H&=L_0+L_1H,&
 A_\eps&=\frac{L_H\Delta}{\eps^2},&
 B_\eps&=\frac{L_1\Delta}{\eps},
 \label{eq:abbrev}\\
 V_\eps&=\frac{\sigma^2}{\eps^2},&
 P_\eps&=1+V_\eps,&
 N_\eps&=1+\min\left\{m,\frac{H^2}{\eps^2}\right\}.
 \label{eq:refined-parameters}
\end{align}
The ratio
\begin{equation}
 t_\eps:=\frac{B_\eps}{A_\eps}=\frac{L_1\eps}{L_H}
 \label{eq:theta}
\end{equation}
separates high accuracy ($t_\eps\le1$) from coarse accuracy
($t_\eps>1$). The parameter $N_\eps$ is an effective component count, while
$P_\eps$ measures fresh-noise cost. Notice that $t_\eps>1$ implies
$H<\eps$ and hence $N_\eps<2$.

Throughout the minimax statements, the nontrivial-accuracy assumption means
\begin{equation}
 0<\eps^2\le c_0\Delta L_H
 \label{eq:nontrivial}
\end{equation}
for one sufficiently small universal constant $c_0>0$. We do not assert a
minimax characterization outside this range. To specify the logarithmic
losses, set
\begin{align}
 \Xi_\eps&=\left[1+\log\left(1+\frac{L_1\eps}{L_0}\right)\right]^4,
 \label{eq:stochastic-coarse-log}\\
 \mathcal L_\eps&=1+\log\left(2+m+A_\eps+B_\eps+P_\eps
                  +\frac H\eps+\frac{L_1\eps}{L_0}\right).
 \label{eq:global-log}
\end{align}
The symbols $\widetilde O$, $\widetilde\Omega$, and
$\widetilde\Theta$ suppress universal powers of these logarithmic factors.
Unadorned $O$, $\Omega$, and $\Theta$ suppress universal constants, unless a
fixed regime constant is explicitly specified.

\subsection{The sample-complexity characterization}

\begin{theorem}[Minimax total sample complexity]
\label{thm:main}
There are universal constants $c_0,c,C>0$ such that, under
\eqref{eq:nontrivial}, for every admissible fixed communication matrix $W$,
\begin{equation}
 \frac{c}{\Xi_\eps}\,\mathfrak F_\eps
 \le S^*\le C\mathcal L_\eps^3\,\mathfrak F_\eps,
 \qquad
 \mathfrak F_\eps:=(A_\eps+B_\eps)P_\eps
                  +A_\eps\sqrt{N_\eps}+N_\eps.
 \label{eq:main-explicit}
\end{equation}
In particular,
\begin{equation}
 \boxed{S^*=\tTheta\!\left(
       (A_\eps+B_\eps)P_\eps
       +A_\eps\sqrt{N_\eps}+N_\eps\right).}
 \label{eq:fully-closed-all-regime-rate}
\end{equation}
The sample upper bound can be attained together with
\begin{equation}
 \E C_{\rm comm}\le C(r(W)-1)(1+A_\eps+B_\eps).
 \label{eq:main-communication}
\end{equation}
The lower bound applies to arbitrary randomized adaptive protocols with
stopping-time query counts and arbitrary output points. In the high-accuracy
region $B_\eps\le A_\eps$, its logarithmic loss can be removed:
$S^*\ge c\mathfrak F_\eps$.
\end{theorem}

The proof is assembled in Section~\ref{sec:assembly}; all constructions and
estimator details are given below. Equation~\eqref{eq:main-communication}
is an achievable communication bound for a sample-efficient method. It is
not a characterization of the joint sample--communication Pareto frontier.

\begin{table}[t]
\centering\small
\caption{Consequences of Theorem~\ref{thm:main}. Every row assumes
\eqref{eq:nontrivial}; logarithmic factors are suppressed.}
\label{tab:regimes}
\renewcommand{\arraystretch}{1.2}
\begin{tabularx}{\textwidth}{@{}p{0.25\textwidth}p{0.39\textwidth}Y@{}}
\toprule
Regime & Minimax total samples & Interpretation\\
\midrule
$B_\eps\le A_\eps$
& $\tTheta(A_\eps P_\eps+A_\eps\sqrt{N_\eps}+N_\eps)$
& Smooth stochastic cost plus component identification.\\
$B_\eps>A_\eps$
& $\tTheta(B_\eps P_\eps)$
& Here $N_\eps<2$; fresh noise multiplies coarse progress.\\
$\sigma=0$, all accuracies
& $\tTheta(A_\eps\sqrt{N_\eps}+N_\eps+B_\eps)$
& Exact local gradients.\\
$H=0$, all accuracies
& $\tTheta((A_\eps+B_\eps)P_\eps)$
& Identical local gradients; $L_H=L_0$ and $N_\eps=1$.\\
\bottomrule
\end{tabularx}
\end{table}

The four lower-bound scales can be proved on separate hard families:
$A_\eps P_\eps$ by an orthogonal smooth embedding,
$N_\eps$ by hidden-sign discovery,
$A_\eps\sqrt{N_\eps}$ by a retiring finite-sum chain, and
$B_\eps P_\eps/\Xi_\eps$ by a terminally flat stochastic chain.
Taking the maximum of four lower bounds loses only a constant relative to
their sum. No direct-sum assertion about simultaneously realizing all four
mechanisms is needed.

\section{Comparison with prior work}
\label{sec:related}

Comparisons require the same counting convention and information interface.
We count \emph{total} local gradient calls. A fresh stochastic response at
one point is different from an exact component gradient, or from a random
sample that can be evaluated at two points. Table~\ref{tab:sample-comparison}
records these distinctions.

\begin{table}[!t]
\centering\small
\caption{Sample-complexity comparisons. Put $A_0=L\Delta/\eps^2$ for an
ordinary smoothness constant $L$. All rates count total gradient calls.
The notation $S_{\rm DNSGD}$ is defined in~\eqref{eq:dnsgd-rounded}.
Each row retains its own assumptions and output guarantee.}
\label{tab:sample-comparison}
\renewcommand{\arraystretch}{1.2}
\begin{tabularx}{\textwidth}{@{}p{0.20\textwidth}p{0.27\textwidth}Y@{}}
\toprule
Work & Rate & Oracle, regime, and guarantee\\
\midrule
Arjevani et al.\ \cite{arjevani2022}; Carmon et al.\ \cite{carmon2019}
& $\Omega(A_0P_\eps)$
& Ordinary $L$-smooth objective; fresh bounded-variance gradients;
$\eps^2\le cL\Delta$; arbitrary randomized algorithms; expected gradient norm at most $\eps$.\\[0.5em]
Zhou et al.\ \cite{snvrg2018}
& $\tO(N_\eps+A_0\sqrt{N_\eps})$
& Exact finite-sum gradients; ordinary averaged squared smoothness $L$;
component variance at most $H^2$; expected squared gradient at most
$\eps^2$.\\[0.5em]
Reisizadeh et al.\ \cite{reisizadeh2023}
& $O(m+A_\eps\sqrt m)$
& Exact finite-sum gradients; averaged relaxed smoothness with constant
$L_H$; $\eps\le cL_H/L_1$; constant-probability stationarity.\\[0.5em]
Peng et al.\ \cite{dense2026}
& $\Omega(m+A_0\sqrt m)$
& Individually $L$-smooth finite sums; exact value--gradient IFO;
$m\ge m_0$, $\eps^2\le cL\Delta$; randomized adaptive queries;
expected-gradient criterion; sufficiently large dimension; no uniform $H$ bound.\\[0.5em]
Luo et al.\ \cite{relaxed2025}
& $O(S_{\rm DNSGD})$
& Local finite-step relaxed smoothness, uniform heterogeneity, and fresh
bounded-variance noise; expected gradient norm at most $\eps$ at every
agent. Integer batches are retained below.\\[0.5em]
This paper
& $\tTheta(\mathfrak F_\eps)$
& Model~\eqref{eq:local-relaxed}--\eqref{eq:output-criterion};
$\eps^2\le c_0\Delta L_H$; expected total calls; dimension unrestricted.
The upper bound works in every finite dimension.\\
\bottomrule
\end{tabularx}
\end{table}

\paragraph{The direct decentralized comparison.}
With the present notation, Corollary~1 of Luo et al.~\cite{relaxed2025}
reports the noise-dependent total cost
$O(V_\eps[A_\eps+1+(B_\eps+1)t_\eps^2])$ after multiplication of their
per-agent formula by $m$. Retaining integer batches gives the useful
uniform upper estimate
\begin{equation}
 S_{\rm DNSGD}:=
 m(1+A_\eps+B_\eps)
 +V_\eps\bigl[A_\eps+1+(B_\eps+1)t_\eps^2\bigr].
 \label{eq:dnsgd-rounded}
\end{equation}
To obtain it directly from their Theorem~1, apply that theorem to
\[
 \bar f_i(y)=\frac{L_H}{\eps^2}
             f_i\!\left(x_0+\frac{\eps}{L_H}y\right).
\]
This makes both the inherited smoothness constant and the target accuracy
one. The quantities $A_\eps,B_\eps,V_\eps,t_\eps$ and the query count
are unchanged. Their choices
then give $T=O(1+A_\eps+B_\eps)$ and
$b=O(1+V_\eps(1+t_\eps^2)/m)$, including rounding and initialization.
The identity $B_\eps=A_\eps t_\eps$ yields
\eqref{eq:dnsgd-rounded}. The corresponding communication guarantee is
$\widetilde O((1+A_\eps+B_\eps)/\sqrt\gamma)$.

For $t_\eps\le1$ and $V_\eps\gtrsim m$, this earlier bound already attains
$O(A_\eps V_\eps)$ under~\eqref{eq:nontrivial}. Our result supplies the
matching necessity of $L_H$ and characterizes the complementary low-noise
and coarse-accuracy regions. In the exact-gradient, full-population regime
$H^2/\eps^2\ge m$, our rate is
$\widetilde\Theta(m+A_\eps\sqrt m)$ at high accuracy. At coarse accuracy,
the fresh-noise term is $\widetilde\Theta(B_\eps P_\eps)$, removing the
extra $t_\eps^2$ in the noise-dependent part of
\eqref{eq:dnsgd-rounded} when that term dominates. These are comparisons
of guarantees under their stated models, rather than claims about the
observed performance of DNSGD.

\paragraph{Variance reduction and component uncertainty.}
The effective population cutoff in $N_\eps$ already appears in nested
variance reduction for ordinarily smooth finite sums~\cite{snvrg2018}.
Variance-reduced clipping gives finite-sum improvements under relaxed
smoothness~\cite{reisizadeh2023}, while normalized methods and SPIDER also
extend to broader generalized-smooth classes~\cite{chen2023}.
Their stochastic $\eps^{-3}$ results use additional mean-squared regularity
and paired evaluations with the same random sample. Our interface supplies
fresh oracle noise, for which the ordinary stochastic obstruction
remains of order $L_H\Delta\sigma^2/\eps^4$~\cite{arjevani2022}.
The dyadic construction separates this cost from component uncertainty.

\paragraph{Lower-bound mechanisms and attribution.}
The inherited smoothness inequality with $L_H$ is Proposition~1 of
Luo et al.~\cite{relaxed2025}. Our orthogonal embedding shows that this
parameter is also necessary for the stated class. The dense sign tables,
bounded radial map, and sequential information principle are adapted from
Peng et al.~\cite{dense2026}; the retiring selector enforces a gradient
heterogeneity bound independent of chain length. The rare-reveal principle
is related to probabilistic zero chains~\cite{arjevani2022}. The coarse
construction adds terminal flatness and a blockwise noncancellation
argument to satisfy the stated finite-step relaxed-smoothness condition.

\paragraph{Communication and quantifiers.}
Finite-time polynomial consensus is classical~\cite{finiteconsensus2007}.
The sample algorithm uses its symmetric known-matrix specialization.
Lu and De Sa~\cite{ludesa2021} establish lower bounds for restricted
synchronous protocols on suitable hard networks, with a specified local
query budget per iteration. Their results do not assert an uncapped
spectral-gap lower bound for every fixed $W$ in the model studied here.
Section~\ref{sec:communication} makes this distinction explicit.

\section{Upper bounds}\label{sec:upper-main}
\subsection{Inherited smoothness and normalized descent}

\begin{lemma}[Inherited smoothness; cf.~\cite{relaxed2025}]
\label{lem:global-relaxed}
Under \eqref{eq:local-relaxed}--\eqref{eq:heterogeneity}, for every
$\norm{x-y}\le1/L_1$,
\begin{equation}
 \norm{\nabla f(x)-\nabla f(y)}
 \le\bigl(L_H+L_1\norm{\nabla f(x)}\bigr)\norm{x-y}.
 \label{eq:global-relaxed}
\end{equation}
Consequently, for every $s$ with $\norm{s}\le1/L_1$,
\begin{equation}
 f(x+s)
 \le f(x)+\inner{\nabla f(x)}{s}
 +\frac12\bigl(L_H+L_1\norm{\nabla f(x)}\bigr)\norm{s}^2.
 \label{eq:descent}
\end{equation}
\end{lemma}

\begin{proof}
Averaging \eqref{eq:local-relaxed} and using
$\norm{\nabla f_i(x)}\le\norm{\nabla f(x)}+H$ gives
\[
\begin{aligned}
 \norm{\nabla f(x)-\nabla f(y)}
 &\le \frac1m\sum_{i=1}^m
       \norm{\nabla f_i(x)-\nabla f_i(y)}\\
 &\le \left(L_0+\frac{L_1}{m}\sum_{i=1}^m
       \norm{\nabla f_i(x)}\right)\norm{x-y}\\
 &\le \bigl(L_0+L_1H+L_1\norm{\nabla f(x)}\bigr)\norm{x-y}.
\end{aligned}
\]
For \eqref{eq:descent}, integrate the directional derivative on the segment
$x+ts$, $t\in[0,1]$, and apply \eqref{eq:global-relaxed} to the pair
$(x,x+ts)$.
\end{proof}

The following perturbation bound relates an inexact normalized direction to
the true gradient.

\begin{lemma}[Direction perturbation]
\label{lem:direction}
Let $g,v\in\R^d$ and define $u=v/\norm v$ when $v\ne0$, and $u=0$ when
$v=0$.  Then
\begin{equation}
 \inner{g}{u}\ge\norm g-2\norm{v-g}.
 \label{eq:direction}
\end{equation}
\end{lemma}

\begin{proof}
If $v\ne0$, then
\[
 \inner g u=\norm v-\inner{v-g}{u}
 \ge\norm v-\norm{v-g}
 \ge\norm g-2\norm{v-g}.
\]
If $v=0$, the right-hand side is $\norm g-2\norm g=-\norm g\le0$.
\end{proof}
\subsection{Exact averaging at a fixed communication matrix}
The next lemma is the symmetric known-spectrum form of finite-time polynomial consensus~\cite{finiteconsensus2007}.

\begin{lemma}[Spectral polynomial consensus]
\label{lem:polynomial-consensus}
Let the distinct eigenvalues of $W$ be
$1=\lambda_1,\lambda_2,\ldots,\lambda_r$.  Define
\begin{equation}
 p_W(t)=\prod_{j=2}^{r}\frac{t-\lambda_j}{1-\lambda_j}.
 \label{eq:consensus-polynomial}
\end{equation}
Then
\begin{equation}
 p_W(W)=J:=\frac1m\one\one^\top.
 \label{eq:exact-consensus}
\end{equation}
Consequently the exact average of any stack of local vectors can be made
available at every agent using $r(W)-1\le m-1$ calls to $W$ and local linear
arithmetic.
\end{lemma}

\begin{proof}
Because $W$ is real symmetric, it has an orthogonal spectral decomposition.
The polynomial satisfies $p_W(1)=1$ and $p_W(\lambda_j)=0$ for every distinct
nonconsensus eigenvalue.  Connectedness makes the eigenspace of $1$ equal to
$\operatorname{span}\{\one\}$, proving \eqref{eq:exact-consensus}.  Horner
or factored evaluation uses one $W$-multiplication per polynomial degree.
\end{proof}

\subsection{A normalized minibatch method}
At a common point $x$, there are two useful constructions.

\paragraph{Random-agent batch.}
Draw public independent indices $I_1,\ldots,I_q\sim\operatorname{Unif}[m]$ and
one fresh local sample for each selected index.  Set
\begin{equation}
 \widehat g_R(x)=\frac1q\sum_{k=1}^q G_{I_k}(x;\xi_k).
 \label{eq:random-agent-estimator}
\end{equation}
It can be formed by letting agent $i$ store
$(m/q)\sum_{k:I_k=i}G_i(x;\xi_k)$ and applying exact averaging once.  With
$\nu^2:=\sigma^2+H^2$,
\begin{equation}
 \E[\widehat g_R(x)]=\nabla f(x),
 \qquad
 \E\norm{\widehat g_R(x)-\nabla f(x)}^2\le\frac{\nu^2}{q}.
 \label{eq:random-agent-variance}
\end{equation}
Indeed, conditional unbiasedness kills the cross term between oracle noise and
$\nabla f_i(x)-\nabla f(x)$.

\paragraph{Full-network batch.}
Every agent takes $b$ fresh samples and exact averaging forms
\begin{equation}
 \widehat g_F(x)=\frac1m\sum_{i=1}^m
 \left(\frac1b\sum_{k=1}^bG_i(x;\xi_{i,k})\right).
 \label{eq:full-network-estimator}
\end{equation}
Then
\begin{equation}
 \E[\widehat g_F(x)]=\nabla f(x),
 \qquad
 \E\norm{\widehat g_F(x)-\nabla f(x)}^2\le\frac{\sigma^2}{mb}.
 \label{eq:full-network-variance}
\end{equation}
Here heterogeneity cancels before the stochastic variance is evaluated.



\begin{theorem}[Normalized minibatch upper bound]
\label{thm:sample-upper}
There is a decentralized algorithm for which all agents return the same
random output $\widehat x$ and
\[
 \E\norm{\nabla f(\widehat x)}\le\eps.
\]
Its total sample and communication costs obey
\begin{align}
 S
 &\le C_1(1+A_\eps+B_\eps)
 \min\left\{
 1+\frac{\sigma^2+H^2}{\eps^2},
 m+\frac{\sigma^2}{\eps^2}
 \right\},
 \label{eq:upper-min-form}\\
 C
 &\le C_2\bigl(r(W)-1\bigr)(1+A_\eps+B_\eps),
 \label{eq:upper-communication-exact}
\end{align}
for universal constants $C_1,C_2$.  In particular,
\begin{equation}
 S=\cO\!\left(
 (1+A_\eps+B_\eps)
 \left[1+\frac{\sigma^2}{\eps^2}
 +\min\left\{m,\frac{H^2}{\eps^2}\right\}
 \right]
 \right).
 \label{eq:upper-simple-form}
\end{equation}
\end{theorem}

\begin{proof}
Choose the cheaper of the following two batch rules:
\begin{align*}
 q_R&=\max\left\{1,
 \left\lceil\frac{256(\sigma^2+H^2)}{\eps^2}\right\rceil\right\},\\
 b_F&=\max\left\{1,
 \left\lceil\frac{256\sigma^2}{m\eps^2}\right\rceil\right\},
 \qquad q_F=mb_F.
\end{align*}
By \eqref{eq:random-agent-variance} or \eqref{eq:full-network-variance}, the
selected estimator $v_t$ satisfies, conditionally on the common iterate $x_t$,
\begin{equation}
 \E[v_t\mid x_t]=g_t:=\nabla f(x_t),
 \qquad
 \E[\norm{v_t-g_t}\mid x_t]\le\frac{\eps}{16}.
 \label{eq:estimator-error}
\end{equation}
Exact consensus from Lemma~\ref{lem:polynomial-consensus} makes $v_t$ available
at every agent.  Define $u_t=v_t/\norm{v_t}$ when $v_t\ne0$ and $u_t=0$
otherwise, and update
\begin{equation}
 x_{t+1}=x_t-\eta u_t,
 \qquad
 \eta=\min\left\{\frac1{2L_1},\frac{\eps}{16L_H}\right\},
 \label{eq:normalized-update}
\end{equation}
where $1/(2L_1)=+\infty$ if $L_1=0$.

Since $\eta\le1/L_1$, Lemmas~\ref{lem:global-relaxed} and
\ref{lem:direction} give
\[
\begin{aligned}
 \E[f(x_{t+1})\mid x_t]
 &\le f(x_t)
 -\eta\E[\inner{g_t}{u_t}\mid x_t]
 +\frac{\eta^2}{2}(L_H+L_1\norm{g_t})\\
 &\le f(x_t)
 -\eta\left(1-\frac{\eta L_1}{2}\right)\norm{g_t}
 +2\eta\E[\norm{v_t-g_t}\mid x_t]
 +\frac{\eta^2L_H}{2}.
\end{aligned}
\]
Let
\[
 T=\max\left\{1,
 \left\lceil\frac{16\Delta}{\eta\eps}\right\rceil\right\}.
\]
Summation, the global gap bound, \eqref{eq:estimator-error},
$1-\eta L_1/2\ge3/4$, and $\eta L_H/2\le\eps/32$ imply
\[
 \frac34\cdot\frac1T\sum_{t=0}^{T-1}
 \E\norm{\nabla f(x_t)}
 \le \frac{\Delta}{\eta T}+\frac{\eps}{8}+\frac{\eps}{32}
 \le\frac{7\eps}{32}.
\]
Thus a public uniform choice from $x_0,\ldots,x_{T-1}$ has expected gradient
norm below $\eps$.

Moreover,
\[
 T=\cO(1+A_\eps+B_\eps),
\]
while
\[
 q_R\le1+\frac{256(\sigma^2+H^2)}{\eps^2},
 \qquad
 q_F\le m+\frac{256\sigma^2}{\eps^2}.
\]
This proves \eqref{eq:upper-min-form}.  Each iteration uses one exact average,
which costs $r(W)-1$ calls, proving \eqref{eq:upper-communication-exact}.
Finally,
\[
 \min\{1+V_\eps+H^2/\eps^2,\,m+V_\eps\}
 =\cO\!\left(1+V_\eps+\min\{m,H^2/\eps^2\}\right),
\]
which gives \eqref{eq:upper-simple-form}.
\end{proof}
\subsection{The refined upper bound}

The minibatch method pays the component uncertainty at every iteration.
The dyadic method reuses reference points while resampling every stochastic
response. It has the following sharper guarantee.

\begin{theorem}[Dyadic fresh-noise upper bound]
\label{thm:dyadic-upper}
There is a decentralized algorithm for which all agents return the same
$\widehat x$ satisfying $\E\norm{\nabla f(\widehat x)}\le\eps$ and
\begin{equation}
 \E S\le\tO\left(
 N_\eps+(1+A_\eps+B_\eps)
                    (\sqrt{N_\eps}+P_\eps)\right).
 \label{eq:dyadic-master-upper}
\end{equation}
Its expected communication cost at the fixed matrix $W$ satisfies
\begin{equation}
 \E C_{\rm comm}\le C(r(W)-1)(1+A_\eps+B_\eps).
 \label{eq:dyadic-communication}
\end{equation}
All local oracle evaluations use fresh samples.
\end{theorem}

\begin{proof}[Proof outline]
Appendix~\ref{sec:dyadic} gives the complete algorithm and proof. It
partitions a virtual path into epochs of length $Q\asymp\sqrt{N_\eps}$
and represents each gradient by an anchor and a binary partition of gradient
differences. A length-$h$ difference costs, up to logarithms,
$\min\{m,h^2,H^2/\eps^2\}+P_\eps$ calls and is formed once per $h$ virtual
steps. Summing over dyadic lengths separates the amortized costs
$\sqrt{N_\eps}$ and $P_\eps$. A finite robust aggregation rule controls
all estimates simultaneously. Normalized descent bounds the expected
virtual path length by $O(1+A_\eps+B_\eps)$.
\end{proof}

\section{Lower bounds and assembly}\label{sec:lower-main}
\subsection{An orthogonal embedding of smooth optimization}

The following reduction realizes the inherited scale $L_H$ using admissible
local functions. It also preserves the noise distribution of any centralized
gradient oracle.

\begin{theorem}[Orthogonal-simplex reduction]
\label{thm:sample-lower}
There are universal constants $c,c_0>0$ such that the following holds.  Let
\[
 0<\eps\le c_0\min\{\sigma,\sqrt{\Delta L_H}\}.
\]
For every randomized adaptive decentralized algorithm in the stated model,
there is a finite-dimensional admissible instance for which success at every
agent requires
\begin{equation}
 \E S\ge c\left(
      \frac{L_H\Delta}{\eps^2}
      +\frac{L_H\sigma^2\Delta}{\eps^4}
      \right).
 \label{eq:sample-lower}
\end{equation}
\end{theorem}

\begin{proof}
Let $g:\R^{d_0}\to\R$ be any ordinary $L_H$-smooth instance with
$g(0)-\inf g\le\Delta$.  In an auxiliary copy of $\R^m$, define
\begin{equation}
 u_i=H\sqrt{\frac{m}{m-1}}
 \left(e_i-\frac1m\one\right),
 \qquad i=1,\ldots,m.
 \label{eq:simplex}
\end{equation}
Then
\[
 \frac1m\sum_{i=1}^m u_i=0,
 \qquad
 \norm{u_i}=H.
\]
For $x=(z,t)\in\R^{d_0}\times\R^m$, set
\begin{equation}
 f_i(z,t)=g(z)+\inner{u_i}{t}.
 \label{eq:embedded-local}
\end{equation}
Their average and gradients are
\begin{equation}
 f(z,t)=g(z),
 \qquad
 \nabla f_i(z,t)=(\nabla g(z),u_i),
 \qquad
 \nabla f(z,t)=(\nabla g(z),0).
 \label{eq:embedded-gradients}
\end{equation}
Thus the global gap is unchanged and the heterogeneity is exactly $H$.
Moreover, for arbitrary $x=(z,t)$ and $y=(z',t')$,
\[
\begin{aligned}
 \norm{\nabla f_i(x)-\nabla f_i(y)}
 &=\norm{\nabla g(z)-\nabla g(z')}\\
 &\le L_H\norm{z-z'}
 \le L_H\norm{x-y}.
\end{aligned}
\]
Since $\norm{\nabla f_i(x)}\ge\norm{u_i}=H$,
\[
 L_H=L_0+L_1H
 \le L_0+L_1\norm{\nabla f_i(x)}.
\]
Therefore every $f_i$ satisfies \eqref{eq:local-relaxed}, in fact globally.

Given any unbiased oracle $G_g(z;\xi)$ for $g$, define the local oracle
\begin{equation}
 G_i((z,t);\xi)=(G_g(z;\xi),u_i).
 \label{eq:embedded-oracle}
\end{equation}
It has exactly the same noise norm and variance as the oracle for $g$.

For a fixed-budget algorithm using $S$ local calls, a centralized
simulator with access to the oracle of $g$ can maintain the complete virtual
state of all agents.  Whenever virtual agent $i$ queries $(z,t)$, the simulator
uses one fresh query at $z$ and returns \eqref{eq:embedded-oracle}.  It can
perform all $W$-multiplications and all local arithmetic internally, without
additional oracle calls.  Hence it reproduces the decentralized transcript
using exactly $S$ queries to $g$.

At termination, the $z$-component of the output of any one agent satisfies
\[
 \E\norm{\nabla g(\widehat z_i)}
 =\E\norm{\nabla f(\widehat z_i,\widehat t_i)}
 \le\eps.
\]
The sharp centralized lower bound for arbitrary randomized adaptive
fresh-sample algorithms therefore yields the fixed-budget form of
\eqref{eq:sample-lower}; see \cite[Theorem~3 and its fixed-budget construction]{arjevani2022}.  For completeness, apply that construction at a fixed universal multiple of
$\eps$, so that the capped output has gradient norm at least $8\eps$ with
probability at least $1/2$. If the worst-case expected number of calls were
a sufficiently small fraction of its fixed-budget threshold, Markov's
inequality would bound the probability of exceeding the cap by $1/8$ on
the same hard instance. The original output would then have expected
gradient norm larger than $\eps$. This proves the expected-call bound.
\end{proof}

\begin{corollary}[The ordinary stochastic scale]
\label{cor:AP-lower}
There are universal constants $c,c_0>0$ such that, whenever
$\eps^2\le c_0\Delta L_H$, the minimax total sample complexity satisfies
\begin{equation}
 S^*\ge cA_\eps P_\eps.
 \label{eq:AP-lower}
\end{equation}
\end{corollary}

\begin{proof}
If $\sigma$ is at least a sufficiently large universal multiple of $\eps$,
Theorem~\ref{thm:sample-lower} applies and gives
$\Omega(A_\eps+A_\eps\sigma^2/\eps^2)=\Omega(A_\eps P_\eps)$.
If $\sigma=O(\eps)$, then $P_\eps=\Theta(1)$, and the classical deterministic
smooth nonconvex lower bound~\cite{carmon2019} gives $\Omega(L_H\Delta/\eps^2)=\Omega(A_\eps)$.
The orthogonal-simplex embedding in the proof of
Theorem~\ref{thm:sample-lower} transfers that deterministic hard instance to
the present decentralized class.  The overlap between the two cases is
absorbed into the universal constants.
\end{proof}

\begin{remark}[Why this proves necessity of $L_1H$]
The local Hessian scale of the hard component is $L_H$, although the primitive
constant in \eqref{eq:local-relaxed} is only $L_0$.  The additional curvature is
certified by the orthogonal gradient component of norm $H$.  Hence the lower
bound contains $L_0+L_1H$ for structural, rather than merely analytic,
reasons.
\end{remark}

\begin{remark}[Scope of the construction]
Our model assumes that the average objective is bounded below; it
does not require every $f_i$ to be bounded below.  The linear orthogonal
components in \eqref{eq:embedded-local} use this scope.  Adding a
separate lower-boundedness requirement on every local objective defines a
strictly different problem and would require another embedding.
\end{remark}
\subsection{Discovering heterogeneous components}
\label{sec:discovery-lower}

The centralized hard instance proves the multiplicative term
$A_\eps P_\eps$, but it does not charge the algorithm for discovering which
local components are present.  The following construction supplies this
additional obstruction and uses the uniform heterogeneity bound exactly.

\begin{theorem}[Hidden-sign discovery lower bound]
\label{thm:discovery-lower}
There are universal constants $c,c_0>0$ such that, whenever
$\eps^2\le c_0\Delta L_H$, every randomized adaptive decentralized algorithm
has an admissible deterministic-oracle instance on which
\begin{equation}
 \E S\ge c\min\left\{m,\frac{H^2}{\eps^2}\right\}.
 \label{eq:discovery-lower}
\end{equation}
Consequently, after combining with the centralized smooth and stochastic
lower bounds,
\begin{equation}
 \boxed{
 S^*\ge c\bigl(A_\eps P_\eps+N_\eps\bigr).}
 \label{eq:combined-sample-lower}
\end{equation}
\end{theorem}

\begin{proof}
Write $K_H=\min\{m,H^2/\eps^2\}$.  If $K_H$ is below a
universal constant, the deterministic smooth lower bound and the nontrivial-
accuracy assumption give $\E S=\Omega(1)=\Omega(K_H)$.  Otherwise it
suffices to prove the claim for every integer
\begin{equation}
 2\le n\le c_1\min\left\{m,\frac{H^2}{\eps^2}\right\}.
 \label{eq:n-discovery-choice}
\end{equation}
Partition the $m$ agents into $n$ nonempty groups $G_1,\ldots,G_n$ with group
sizes $g_j$ differing by at most one, and set
\[
 w_j:=\frac{m}{ng_j}.
\]
Then $1/2\le w_j\le2$.  Draw independent hidden signs
$s_1,\ldots,s_n\in\{-1,+1\}$, let
\[
 L_\star:=L_0+\frac12L_1H,
 \qquad
 a:=4\eps\sqrt n,
\]
and, for an agent $i\in G_j$, define on $z\in\R^n$
\begin{equation}
 h_i(z)=\frac{L_\star}{2}\norm z^2-a w_j s_jz_j.
 \label{eq:sign-local}
\end{equation}
The average is
\begin{equation}
 h(z)=\frac{L_\star}{2}\norm z^2
 -\frac an\sum_{j=1}^ns_jz_j.
 \label{eq:sign-average}
\end{equation}

Append an auxiliary orthogonal coordinate and choose simplex vectors $u_i$
of norm $H/2$ and mean zero.  Set
\begin{equation}
 f_i(z,t)=h_i(z)+\inner{u_i}{t}.
 \label{eq:sign-embedded}
\end{equation}
The base-gradient disagreement is bounded by
\begin{equation}
 \norm{a w_js_je_j-\frac an(s_1,\ldots,s_n)}
 \le\sqrt5\,a.
 \label{eq:sign-base-heterogeneity}
\end{equation}
Taking $c_1$ sufficiently small makes the squared sum of this quantity and
$(H/2)^2$ at most $H^2$.  Thus \eqref{eq:heterogeneity} holds.  Every local
objective is ordinarily $L_\star$-smooth and
$\norm{\nabla f_i}\ge H/2$, so
\[
 L_\star=L_0+\frac12L_1H
 \le L_0+L_1\norm{\nabla f_i}.
\]
Hence the relaxed-smoothness requirement also holds globally.

At the common start $z_0=0$,
\begin{equation}
 h(0)-\inf h
 =\frac{a^2}{2nL_\star}
 =\frac{8\eps^2}{L_\star}
 \le\Delta
 \label{eq:sign-gap}
\end{equation}
under the stated nontrivial-accuracy condition, after adjusting the universal
constant.

One query to a member of $G_j$ can reveal $s_j$, but no query outside $G_j$
contains that sign.  Let $D$ be the number of distinct groups queried before
termination.  Conditional on the complete transcript, every sign in an
unqueried group remains independent and symmetric.  If $U$ denotes the
unqueried groups, then for the transcript-measurable output $z$ of any fixed
agent,
\[
 \nabla_U h(z)=L_\star z_U-\frac an s_U.
\]
The function
\[
 z_U\longmapsto
 \E_{s_U}\norm{L_\star z_U-\frac an s_U}
\]
is convex and even, hence minimized at $z_U=0$.  On $D<n/2$,
\begin{equation}
 \E\bigl[\norm{\nabla h(z)}\mid\text{transcript}\bigr]
 \ge\frac an\sqrt{|U|}
 \ge\frac{a}{\sqrt{2n}}
 =2\sqrt2\,\eps.
 \label{eq:sign-conditional-lower}
\end{equation}
Suppose an algorithm had worst-case expected total call count at most $n/4$.
Under the uniform sign prior, $D\le S$, so Markov's inequality gives
\[
 \Pr(D\ge n/2)\le\frac{2\E S}{n}\le\frac12.
\]
Averaging \eqref{eq:sign-conditional-lower} over the transcript yields
$\E\norm{\nabla h(\widehat z)}>\eps$.  Hence some fixed sign realization
violates the required expected stationarity guarantee.  Thus the worst-case
expected number of calls is $\Omega(n)$.  Maximizing $n$ subject to
\eqref{eq:n-discovery-choice} proves \eqref{eq:discovery-lower}.

For \eqref{eq:combined-sample-lower},
Corollary~\ref{cor:AP-lower} gives $\Omega(A_\eps P_\eps)$.  Under the
nontrivial-accuracy assumption this term is also $\Omega(1)$, so it absorbs
the additive one in $N_\eps=1+K_H$.  The maximum of the two lower bounds is
within a universal factor of their sum.
\end{proof}
\subsection{Sequential component and coarse-noise lower bounds}

The component-discovery construction is a single estimation task. The next
result shows that component uncertainty can recur throughout an optimization
trajectory, even under a uniform heterogeneity bound.

\begin{theorem}[Bounded-disagreement dense-weak-hiding lower bound]
\label{thm:retiring-dwh-lower}
There are universal constants $c,c_0>0$ such that, whenever
\begin{equation}
 \eps^2\le c_0\Delta L_H,
 \label{eq:retiring-nontrivial}
\end{equation}
every randomized adaptive decentralized algorithm has an admissible
exact-gradient instance for which
\begin{equation}
 \boxed{
 \E S\ge c A_\eps\sqrt{N_\eps}.}
 \label{eq:retiring-square-root-lower}
\end{equation}
The algorithm may adaptively choose both the queried agent and an arbitrary
query point from its complete preceding history.
\end{theorem}

\begin{proof}[Construction and proof location]
Appendix~\ref{subsec:retiring-dwh} constructs a chain whose component-specific
signal disappears as a gate opens. A dimension-independent selector bound
keeps the disagreement at $O(\eps/\rho)$ after scaling. Choosing
$\rho\asymp\max\{m^{-1/2},\eps/H\}$ yields
$\Theta(A_\eps\rho)$ sequential stages, each requiring
$\Omega(\min\{m,\rho^{-2}\})$ component queries. Their product is
$\Omega(A_\eps\sqrt{N_\eps})$; the boundary cases follow from the
ordinary and discovery lower bounds.
\end{proof}

At coarse accuracy, the question is whether the path-length scale
$B_\eps$ also incurs the fresh-noise factor $P_\eps$. A hidden stage must
have small gradient amplitude for the variance budget to conceal it for
$\Theta(P_\eps)$ calls. The following construction achieves this while
retaining the prescribed relaxed smoothness.

\begin{theorem}[Stochastic coarse product lower bound]
\label{thm:stochastic-coarse-product}
There are universal constants $c,C>0$ such that the following holds.  Assume
\begin{equation}
 L_1\eps\ge2L_0,
 \qquad
 \eps^2\le c\Delta L_H.
 \label{eq:stochastic-coarse-premise}
\end{equation}
With $\Xi_\eps$ as in~\eqref{eq:stochastic-coarse-log},
assume in addition that $B_\eps\ge C\Xi_\eps$.  Then every randomized
adaptive decentralized algorithm has an admissible
instance with identical local objectives and independent fresh local oracle
samples for which
\begin{equation}
 \boxed{
 \E S\ge
 c\frac{B_\eps P_\eps}{\Xi_\eps}.}
 \label{eq:BP-product-lower}
\end{equation}
The construction has zero heterogeneity and therefore belongs to the class
for every prescribed upper bound $H\ge0$.
\end{theorem}

\begin{proof}[Construction and proof location]
Appendix~\ref{sec:stochastic-coarse-closure} constructs
$T=\Theta(B_\eps/\Xi_\eps)$ hidden blocks with an $8\eps$ stationarity
obstruction. The hidden part of every reply has norm $O(\eps)$, so it can
be revealed with probability $\Theta(P_\eps^{-1})$ under the variance
budget. A prefix-truncated coupling prevents arbitrary adaptive queries or
the final output from bypassing the hidden blocks. The required number of
calls is therefore $\Omega(TP_\eps)$. The appendix proves every analytic
and oracle property, and includes the stopping-time argument.
\end{proof}

\subsection{Proof of the main theorem}
\label{sec:assembly}

\begin{proof}[Proof of Theorem~\ref{thm:main}]
Corollary~\ref{cor:AP-lower} and Theorems~\ref{thm:discovery-lower}
and~\ref{thm:retiring-dwh-lower} give the three lower bounds
\[
 \Omega(A_\eps P_\eps),\qquad \Omega(N_\eps),\qquad
 \Omega(A_\eps\sqrt{N_\eps}).
\]
Here the additive one in $N_\eps$ is absorbed by $A_\eps P_\eps$, since
\eqref{eq:nontrivial} implies $A_\eps\ge1/c_0$.
If $B_\eps\le A_\eps$, these bounds already give
$S^*\ge c\mathfrak F_\eps$.

For the coarse term, if $L_1\eps<2L_0$, then $B_\eps<2A_\eps$ and the
ordinary stochastic bound suffices. If $L_1\eps\ge2L_0$ and
$B_\eps\ge C\Xi_\eps$, Theorem~\ref{thm:stochastic-coarse-product} gives
$S^*\ge cB_\eps P_\eps/\Xi_\eps$. In the remaining case,
$B_\eps<C\Xi_\eps$, and $A_\eps\ge1/c_0$ implies
$B_\eps P_\eps/\Xi_\eps\le Cc_0A_\eps P_\eps$.
Thus all cases give $S^*\ge c\mathfrak F_\eps/\Xi_\eps$ after combining
the four lower bounds. Every construction uses a finite dimension, and the
centralized simulations apply to any fixed admissible $W$.

For the upper bound, Appendix~\ref{sec:dyadic} gives the quantitative estimate
\[
 \E S\le C\mathcal L_\eps^3
 \left[N_\eps+P_\eps+
       (1+A_\eps+B_\eps)(\sqrt{N_\eps}+P_\eps)\right].
\]
If $B_\eps\le A_\eps$, every term containing $B_\eps$ is absorbed by the
corresponding $A_\eps$ term. If $B_\eps>A_\eps$, then $N_\eps<2$, so
$B_\eps\sqrt{N_\eps}\le\sqrt2B_\eps P_\eps$.
The remaining startup terms are absorbed using $A_\eps\ge1/c_0$ and
$\sqrt{N_\eps}\le N_\eps$. This proves the upper inequality in
\eqref{eq:main-explicit}. The same implementation uses
$O(1+A_\eps+B_\eps)$ expected exact averaging operations, each costing
$r(W)-1$ calls, which proves~\eqref{eq:main-communication}.
\end{proof}

\section{Communication at a fixed matrix}
\label{sec:communication}

The sample characterization does not determine the minimum number of
communication rounds. The latter depends on the protocol, the matrix, and
whether sample cost is constrained. We give two explicit implications of
the model used in this paper.

\subsection{Finite spectral degree and the spectral gap}

\begin{proposition}[Two-agent spectral-gap counterexample]
\label{prop:two-agent}
For every $0<\delta\le1/2$, let
\begin{equation}
 W_\delta=
 \begin{pmatrix}
 1-\delta&\delta\\
 \delta&1-\delta
 \end{pmatrix}.
 \label{eq:two-agent-W}
\end{equation}
Then $W_\delta$ is symmetric, doubly stochastic and positive semidefinite, and
its spectral gap is $\gamma=2\delta$.  Nevertheless, one call to $W_\delta$
lets each agent recover the other agent's entire transmitted vector exactly.
Consequently no lower bound containing an uncapped factor
$\gamma^{-1/2}$ can hold pointwise for all fixed $W$ under the stated
communication oracle.
\end{proposition}

\begin{proof}
If the transmitted vectors are $v_1,v_2$, agent $1$ receives
\[
 y_1=(1-\delta)v_1+\delta v_2.
\]
It knows $v_1$ and $\delta$, hence computes
\[
 v_2=\frac{y_1-(1-\delta)v_1}{\delta}.
\]
Agent $2$ argues symmetrically.  This remains true as $\delta\downarrow0$,
while $\gamma^{-1/2}\to\infty$.
\end{proof}


Lemma~\ref{lem:polynomial-consensus} gives the corresponding general
statement: exact averaging costs at most $r(W)-1\le m-1$ calls. The
coefficients may have either sign and may be large; this is permitted by
the exact-real model. A finite-precision implementation would require a
separate numerical-stability analysis.

Accelerated approximate consensus instead gives a useful dependence on
$\gamma^{-1/2}$ and on a logarithmic accuracy parameter~\cite{relaxed2025}.
These two guarantees concern different implementations. The two-agent
example shows that the spectral gap alone does not characterize pointwise
communication complexity when exact polynomial arithmetic is available.

\begin{remark}[Iteration bounds and communication bounds]
A proof that allows a bounded number of local queries per iteration must
retain that restriction when its iteration lower bound is interpreted as
a communication lower bound. Unrestricted local work does not eliminate
all communication obstructions, but it changes which arguments apply.
The network-family and protocol quantifiers in~\cite{ludesa2021} are
therefore essential; Proposition~\ref{prop:two-agent} does not contradict
those results.
\end{remark}

\subsection{A communication-free regime}

\begin{theorem}[Small heterogeneity permits zero communication]
\label{thm:zero-communication}
For any $\sigma\ge0$, if $H\le\eps/32$, each agent can independently
return an output satisfying~\eqref{eq:output-criterion}, with
\begin{equation}
 C_{\rm comm}=0,\qquad
 S=O\bigl(m(1+A_\eps+B_\eps)P_\eps\bigr).
 \label{eq:zero-communication-rate}
\end{equation}
\end{theorem}

\begin{proof}
At its own point $x_{i,t}$, agent $i$ averages
$b=\max\{1,\lceil256\sigma^2/\eps^2\rceil\}$ fresh responses to obtain
$v_{i,t}$. Conditional on that point,
\[
 \E\norm{v_{i,t}-\nabla f(x_{i,t})}
 \le H+\frac{\eps}{16}\le\frac{3\eps}{32}.
\]
It takes a normalized step of length
$\eta=\min\{1/(2L_1),\eps/(16L_H)\}$, with the zero convention.
The direction and descent inequalities give
\[
 \E f(x_{i,t+1})\le\E f(x_{i,t})
 -\frac34\eta\E\norm{\nabla f(x_{i,t})}
 +\frac{3\eta\eps}{16}+\frac{\eta^2L_H}{2}.
\]
Taking $T=\max\{1,\lceil16\Delta/(\eta\eps)\rceil\}$ and summing yields
\[
 \frac34\,\frac1T\sum_{t<T}\E\norm{\nabla f(x_{i,t})}
 \le\frac{\eps}{16}+\frac{3\eps}{16}+\frac{\eps}{32}
 =\frac{9\eps}{32}.
\]
A uniform local iterate therefore meets the target. Since
$T=O(1+A_\eps+B_\eps)$ and $b=O(P_\eps)$, summing over the $m$ agents
proves~\eqref{eq:zero-communication-rate}.
\end{proof}

For $H\le\eps/32$, the main theorem has sample scale
$(A_\eps+B_\eps)P_\eps$ up to logarithmic factors.
Theorem~\ref{thm:zero-communication} provides a communication-free
alternative with an $m$-fold total sample budget. It does not determine the
full sample--communication frontier. In particular, a positive
communication lower bound uniform in $H/\eps$ cannot hold for this class.

\section{Discussion}

The sample rate separates four costs that are conflated by a single
minibatch analysis. Heterogeneity affects both the inherited curvature
$L_H$ and the effective population $N_\eps$, but fresh within-component
noise enters through its own factor $P_\eps$. At coarse accuracy, a
terminally flat chain forces that factor to multiply $B_\eps$ as well.
The explicit logarithmic losses in~\eqref{eq:main-explicit} locate the
remaining difference between the proved upper and lower estimates.

Several questions remain outside the theorem. The lower constructions do
not determine the minimax rate in every fixed ambient dimension. The
one-dimensional construction in Appendix~\ref{sec:coarse-reset} proves a
separate deterministic coarse obstruction, but does not settle that broader
question. Requiring every local objective to be bounded below would also
change the class used by the simplex reductions. Finally, a complete joint
sample--communication frontier would need to track the matrix and allowed
local work explicitly. The bounds and counterexamples in
Section~\ref{sec:communication} identify constraints on such a theory.

\clearpage
\begin{thebibliography}{99}

\bibitem{relaxed2025}
Luo Luo, Xue Cui, Tingkai Jia, and Cheng Chen.
\newblock Decentralized stochastic nonconvex optimization under the
$(L_0,L_1)$-smoothness.
\newblock In \emph{Proceedings of the 32nd ACM SIGKDD Conference on Knowledge
Discovery and Data Mining V.2}, 2026.
\newblock \href{https://doi.org/10.1145/3770855.3817843}{doi:10.1145/3770855.3817843}.
\newblock Also arXiv:2509.08726v3.

\bibitem{ludesa2021}
Yucheng Lu and Christopher De Sa.
\newblock Optimal complexity in decentralized training.
\newblock In \emph{Proceedings of the 38th International Conference on
Machine Learning}, volume 139 of \emph{Proceedings of Machine Learning
Research}, pages 7111--7123, 2021.
\newblock \url{https://proceedings.mlr.press/v139/lu21a.html}.

\bibitem{chen2023}
Ziyi Chen, Yi Zhou, Yingbin Liang, and Zhaosong Lu.
\newblock Generalized-smooth nonconvex optimization is as efficient as smooth
nonconvex optimization.
\newblock In \emph{Proceedings of the 40th International Conference on
Machine Learning}, volume 202 of \emph{Proceedings of Machine Learning
Research}, pages 5396--5427, 2023.
\newblock \url{https://proceedings.mlr.press/v202/chen23ar.html}.

\bibitem{snvrg2018}
Dongruo Zhou, Pan Xu, and Quanquan Gu.
\newblock Stochastic nested variance reduction for nonconvex optimization.
\newblock In \emph{Advances in Neural Information Processing Systems},
volume 31, 2018.
\newblock \url{https://arxiv.org/abs/1806.07811}.

\bibitem{reisizadeh2023}
Amirhossein Reisizadeh, Haochuan Li, Subhro Das, and Ali Jadbabaie.
\newblock Variance-reduced clipping for non-convex optimization.
\newblock In \emph{IEEE International Conference on Acoustics, Speech and
Signal Processing (ICASSP)}, pages 1--5, 2025.
\newblock Full version: arXiv:2303.00883v2, 2023.
\newblock \url{https://arxiv.org/abs/2303.00883}.

\bibitem{dense2026}
Yuxing Peng, Zhiqing Tang, and Weijia Jia.
\newblock Dense weak hiding: Closing complexity gaps in nonconvex and PL
finite-sum optimization under individual smoothness.
\newblock arXiv:2609.00045v3, 2026.
\newblock \url{https://arxiv.org/abs/2609.00045v3}.

\bibitem{arjevani2022}
Yossi Arjevani, Yair Carmon, John C. Duchi, Dylan J. Foster, Nathan Srebro,
and Blake Woodworth.
\newblock Lower bounds for non-convex stochastic optimization.
\newblock \emph{Mathematical Programming}, 199(1--2):165--214, 2023.
\newblock Full version: arXiv:1912.02365v2.
\newblock \url{https://arxiv.org/abs/1912.02365}.

\bibitem{carmon2019}
Yair Carmon, John C. Duchi, Oliver Hinder, and Aaron Sidford.
\newblock Lower bounds for finding stationary points I.
\newblock \emph{Mathematical Programming}, 184(1--2):71--120, 2020.
\newblock \url{https://arxiv.org/abs/1710.11606}.

\bibitem{finiteconsensus2007}
Shreyas Sundaram and Christoforos N. Hadjicostis.
\newblock Finite-time distributed consensus in graphs with time-invariant
 topologies.
\newblock In \emph{American Control Conference}, pages 711--716, 2007.
\newblock \href{https://doi.org/10.1109/ACC.2007.4282726}{doi:10.1109/ACC.2007.4282726}.

\end{thebibliography}

\clearpage
\appendix
\section{A robust dyadic estimator with fresh oracle noise}
\label{sec:dyadic}

The distinction between component sampling and fresh oracle noise is useful
only if it is preserved under adaptive choices of the iterates.  We use a
binary family of reference intervals and a robust estimate for each interval.
The resulting guarantee holds simultaneously for all constructed objects.
This avoids applying an unconditional second-moment estimate at a random
virtual time.

\subsection{Averaged component variation}
\label{sec:averaged-smoothness}

The multilevel estimator below needs a deterministic bound on the average
variation of the local gradients.  No samplewise coupling of the stochastic
oracle is assumed.

\begin{lemma}[Averaged relaxed smoothness]
\label{lem:averaged-relaxed}
For every $x,y$ with $\norm{x-y}\le1/L_1$,
\begin{equation}
 \left(
 \frac1m\sum_{i=1}^m
 \norm{\nabla f_i(x)-\nabla f_i(y)}^2
 \right)^{1/2}
 \le
 \bigl(L_H+L_1\norm{\nabla f(x)}\bigr)\norm{x-y}.
 \label{eq:averaged-relaxed}
\end{equation}
\end{lemma}

\begin{proof}
For every $i$,
\[
 L_0+L_1\norm{\nabla f_i(x)}
 \le L_0+L_1H+L_1\norm{\nabla f(x)}
 =L_H+L_1\norm{\nabla f(x)}.
\]
Apply the local relaxed-smoothness assumption componentwise and take the root
mean square.
\end{proof}

\subsection{Fresh finite-population estimates}

\begin{lemma}[Fresh anchor and difference estimates]
\label{lem:fresh-finite-pop}
Let $a,b$ be measurable with respect to the history before the next oracle
calls.  Draw a uniform subset ${\cal S}\subseteq[m]$ of size $c$ without
replacement.  For each selected component, average $r$ fresh responses at
each required endpoint.  Then
\[
 \widehat D_{a,b}=\frac1c\sum_{i\in{\cal S}}
       \bigl(\bar G_i(b)-\bar G_i(a)\bigr)
\]
is conditionally unbiased for $D_{a,b}=\nabla f(b)-\nabla f(a)$ and satisfies
\begin{equation}
 \E\bigl[\norm{\widehat D_{a,b}-D_{a,b}}^2\mid\mathcal F\bigr]
 \le \frac{\mathbf1_{\{c<m\}}}{c}\,
       \frac1m\sum_{i=1}^m\norm{D^i_{a,b}-D_{a,b}}^2
       +\frac{2\sigma^2}{cr},
 \label{eq:difference-primitive}
\end{equation}
where $D^i_{a,b}=\nabla f_i(b)-\nabla f_i(a)$.
The corresponding anchor $\widehat g(a)=c^{-1}\sum_{i\in{\cal S}}\bar G_i(a)$
is conditionally unbiased and satisfies
\begin{equation}
 \E\bigl[\norm{\widehat g(a)-\nabla f(a)}^2\mid\mathcal F\bigr]
 \le \frac{\mathbf1_{\{c<m\}}H^2}{c}+\frac{\sigma^2}{cr}.
 \label{eq:anchor-primitive}
\end{equation}
\end{lemma}

\begin{proof}
The finite-population variance formula bounds the component-sampling term.
For $c=m$ that term vanishes.  The fresh oracle errors are conditionally
centered martingale differences, so their cross terms have zero expectation.
Their variances sum to the displayed noise terms.  No common random sample is
used at two points.
\end{proof}

\begin{lemma}[A finite robust aggregation rule]
\label{lem:robust-center}
Let $Y_1,\ldots,Y_k$ estimate the same vector $\mu$, where $k$ is odd.
For each $j$, let $r_j$ be the $(k+1)/2$-th smallest member of
$\{\norm{Y_j-Y_i}:1\le i\le k\}$, and return
$\mathsf M(Y_1,\ldots,Y_k)=Y_{j_*}$ for the smallest index minimizing $r_j$.
If, conditionally on the history preceding each replicate,
\[
 \E\norm{Y_j-\mu}^2\le V,
\]
then
\begin{equation}
 \E\norm{\mathsf M(Y_1,\ldots,Y_k)-\mu}^2\le18V.
 \label{eq:robust-global-mse}
\end{equation}
If additionally $V\le a^2/16$, then
\begin{equation}
 \Pr\{\norm{\mathsf M(Y_1,\ldots,Y_k)-\mu}>3a\}
 \le2^{-k}.
 \label{eq:robust-tail}
\end{equation}
The conclusions also hold conditionally on the history preceding all
replicates.  Independence beyond the stated conditional moment bounds is
unnecessary.
\end{lemma}

\begin{proof}
Write $q=(k+1)/2$.  If at least $q$ replicates lie within distance $a$ of
$\mu$, each of these replicates has $q$ neighbors within distance $2a$.
Consequently $r_{j_*}\le2a$.  The selected ball and the set of $q$ accurate
replicates intersect, so the selected center lies within $3a$ of $\mu$.
More generally, if $R$ is the $q$-th smallest error
$\norm{Y_j-\mu}$, the same argument gives
$\norm{\mathsf M-\mu}\le3R$.  Since at least $q$ errors are at least $R$,
\[
 \norm{\mathsf M-\mu}^2
 \le \frac{18}{k}\sum_{j=1}^k\norm{Y_j-\mu}^2.
\]
Taking expectations proves \eqref{eq:robust-global-mse}.
For \eqref{eq:robust-tail}, conditional Markov inequalities bound the
probability of each replicate error exceeding $a$ by $1/16$.  Iterated
conditioning bounds the probability that any specified $q$ replicates are
inaccurate by $16^{-q}$.  A union bound over subsets gives
$2^k16^{-q}\le2^{-k}$.
\end{proof}

\subsection{Normalized steps and a bounded virtual path}

Put $L=L_H$, set $\delta=\eps/32$, and choose
\begin{equation}
 \eta=\min\left\{\frac1{4L_1},\frac{\delta}{4L}\right\},
 \qquad
 T=\max\left\{1,\left\lceil\frac{16\Delta}{\eta\delta}\right\rceil\right\}.
 \label{eq:dyadic-macro-parameters}
\end{equation}
At macro-iterate $x_t$, the algorithm uses an estimate $v_t$ and takes the
normalized step $x_{t+1}=x_t-\eta u_t$, with
$u_t=v_t/\norm{v_t}$ and the usual zero convention.  To construct estimates
at future points, subdivide this segment into
\begin{equation}
 J_t=\left\lceil8\left(1+
                 \frac{\eta L_1\norm{v_t}}{\delta}\right)\right\rceil,
 \qquad
 z_{t,j}=x_t-\frac{j\eta}{J_t}u_t,
 \quad 0\le j\le J_t.
 \label{eq:intrinsic-micro-path}
\end{equation}
These intermediate points affect only the estimator; they do not introduce
additional descent decisions.

We impose a deterministic safeguard to control the contribution of rare
estimation failures.  Define
\begin{equation}
 \mathcal G(D)=2\sqrt{LD}+4L_1D,\qquad
 \overline\Delta=\Delta+3T\eta\delta,\qquad
 M=2\bigl(\mathcal G(\overline\Delta)+\delta\bigr).
 \label{eq:dyadic-safeguard}
\end{equation}
If $\norm{v_t}>M$, terminate and return $x_0$.  Otherwise accept the
macro-iterate and continue.  Consequently
\begin{equation}
 J_t\le J_{\max}:=
 \left\lceil8\left(1+\frac{\eta L_1M}{\delta}\right)\right\rceil,
 \qquad U\le U_{\max}:=TJ_{\max}.
 \label{eq:dyadic-deterministic-cap}
\end{equation}
If no safeguard is triggered, return a fresh public uniform choice from
$x_0,\ldots,x_{T-1}$.

\begin{lemma}[Gap and pathwise variation bounds]
\label{lem:micro-variation}
At every point $x$ with $f(x)-\inf f\le D$,
\begin{equation}
 \norm{\nabla f(x)}\le\mathcal G(D).
 \label{eq:gap-gradient-envelope}
\end{equation}
If $\norm{v_t-\nabla f(x_t)}\le\delta$, each micro-edge of
\eqref{eq:intrinsic-micro-path} satisfies
\begin{equation}
 \max_i\norm{\nabla f_i(z_{t,j+1})-
                      \nabla f_i(z_{t,j})}\le\delta.
 \label{eq:micro-edge-variation}
\end{equation}
Thus, on any path whose preceding macro-estimates satisfy this error bound,
a block $I$ of $h$ consecutive micro-edges obeys
\begin{equation}
 \frac1m\sum_i\norm{D^i_I-D_I}^2
 \le\min\{h^2\delta^2,4H^2\}.
 \label{eq:block-variation}
\end{equation}
\end{lemma}

\begin{proof}
Let $g=\norm{\nabla f(x)}$.  For $g>0$, apply the descent inequality to a
step of length $s=\min\{g/(2L),1/(2L_1)\}$ in the negative gradient
direction.  Since $(L+L_1g)s\le g$, the objective decreases by at least
$sg/2$.  It follows that
$D\ge\min\{g^2/(4L),g/(4L_1)\}$, which proves
\eqref{eq:gap-gradient-envelope}.

Every point $z$ on an accepted macro-segment satisfies
\[
 \norm{\nabla f(z)}
 \le(1+\eta L_1)\norm{\nabla f(x_t)}+\eta L
 \le\frac54\norm{v_t}+\frac32\delta.
\]
Apply the local smoothness inequality to a micro-edge and use
$\norm{\nabla f_i(z)}\le\norm{\nabla f(z)}+H$.  The edge variation is at most
\[
 \frac{\eta}{J_t}
 \left(L+\frac54L_1\norm{v_t}+\frac32L_1\delta\right)
 \le\delta,
\]
by \eqref{eq:dyadic-macro-parameters}--\eqref{eq:intrinsic-micro-path}.
Telescoping gives $\norm{D^i_I}\le h\delta$, and centering cannot increase
the average square.  The alternative bound follows from
$\norm{D^i_I-D_I}\le2H$.
\end{proof}

\subsection{Binary reference intervals and simultaneous accuracy}

Choose a power of two $Q$ and an integer $\ell$ such that
\begin{equation}
 Q\le\sqrt{N_\eps}<2Q,\qquad \ell=2+\log_2Q.
 \label{eq:dyadic-Q}
\end{equation}
Partition the virtual path into epochs of $Q$ edges.  At each epoch start,
form a fresh anchor.  Within an epoch, maintain the binary partition of the
prefix: append a unit interval, and whenever two adjacent active intervals
have the same length, replace them by a fresh estimate across their union.
At every virtual time there are at most $\ell$ active objects, including the
anchor.  All endpoints are fixed before the corresponding fresh calls.

Fix $C_b=2^{16}$.  For an anchor, use
\begin{equation}
 c_0=\min\left\{m,\max\left\{1,
       \left\lceil C_b\ell^2H^2/\delta^2\right\rceil\right\}\right\},
 \qquad
 r_0=\max\left\{1,
       \left\lceil C_b\ell^2\sigma^2/(c_0\delta^2)\right\rceil\right\}.
 \label{eq:dyadic-anchor-batches}
\end{equation}
For an interval of length $h$, use
\begin{equation}
 \begin{split}
 c_h&=\min\left\{m,\max\left\{1,
       \left\lceil C_b\ell^2\min\{h^2,H^2/\delta^2\}\right\rceil
                              \right\}\right\},\\
 r_h&=\max\left\{1,
       \left\lceil C_b\ell^2\sigma^2/(c_h\delta^2)\right\rceil\right\}.
 \end{split}
 \label{eq:dyadic-block-batches}
\end{equation}
For each object, construct $k$ fresh replicate estimates from
Lemma~\ref{lem:fresh-finite-pop} and aggregate them by
Lemma~\ref{lem:robust-center}.  In particular, the aggregate need not be
unbiased.  The subsets and all endpoint responses are drawn afresh in every
replicate.

Here are explicit choices of the confidence parameters.  Set
\begin{equation}
 M_{\rm obj}=(\ell+2)(U_{\max}+1),\qquad
 V_*=128(H^2+\delta^2),
 \label{eq:dyadic-object-envelope}
\end{equation}
and
\begin{equation}
 p=\min\left\{\frac14,\frac1{U_{\max}+1},
       \frac{\delta}{4(M+\delta)},
       \frac{\delta^2}{16\ell M_{\rm obj}V_*}\right\}.
 \label{eq:dyadic-failure-budget}
\end{equation}
Take $k$ to be the smallest odd integer at least
$\log_2(M_{\rm obj}/p)$.  These choices introduce only logarithmic factors
in the dimensionless problem parameters.

At virtual time $u$, let ${\cal P}(u)$ denote the active intervals and define
\begin{equation}
 w_u=\widehat g(z_{u_0})+
           \sum_{I\in{\cal P}(u)}\widehat D_I,
 \label{eq:dyadic-estimator}
\end{equation}
where $u_0$ is the epoch start.  The macro-estimator $v_t$ is $w_u$ at the
virtual time corresponding to $x_t$.

\begin{lemma}[Simultaneous accuracy and cost]
\label{lem:dyadic-invariant}
With probability at least $1-p$, every constructed anchor or difference has
error at most $\delta/(2\ell)$.  On this event, every virtual estimate has
error at most $\delta/2$, including those at random macro endpoints, and the
safeguard never triggers.  Moreover,
\begin{equation}
 \E S\le\tO\left(N_\eps+P_\eps+
            \E U\bigl(\sqrt{N_\eps}+P_\eps\bigr)\right).
 \label{eq:dyadic-path-cost}
\end{equation}
Even without restricting to the simultaneous-accuracy event, each
constructed object's error has conditional second moment at most $V_*$.
\end{lemma}

\begin{proof}
Suppose all previously constructed objects are accurate.  The true gradient
telescopes over the active partition, so each preceding macro-estimate has
error at most $\delta/2$.  Lemma~\ref{lem:micro-variation} then supplies the
pathwise variation bound for the next object's endpoints.  Equations
\eqref{eq:difference-primitive}--\eqref{eq:dyadic-block-batches} bound the
conditional variance of every fresh replicate by
\[
 \frac{6\delta^2}{C_b\ell^2}
 \le\frac1{16}\left(\frac{\delta}{6\ell}\right)^2.
\]
Lemma~\ref{lem:robust-center}, with $a=\delta/(6\ell)$, bounds the
conditional probability that the next aggregate error exceeds
$\delta/(2\ell)$ by $2^{-k}$.  This is a first-failure argument: no
conditioning on future accuracy events is used.

At most $M_{\rm obj}$ objects can be created before termination, by
\eqref{eq:dyadic-deterministic-cap}.  A union bound therefore gives a total
failure probability at most $M_{\rm obj}2^{-k}\le p$.
On the resulting event, summing the at most $\ell$ active errors gives the
claimed pathwise estimate, irrespective of the random macro endpoint.
The normalized descent inequality also gives
\[
 f(x_{t+1})\le f(x_t)+3\eta\delta.
\]
Hence all accepted iterates have gap at most $\overline\Delta$.
Equation~\eqref{eq:gap-gradient-envelope} and the estimator error bound imply
$\norm{v_t}\le\mathcal G(\overline\Delta)+\delta<M$.
Thus the safeguard cannot trigger on this event.

For the unconditional moment claim, uniform heterogeneity always gives
\[
 m^{-1}\sum_i\norm{D^i_I-D_I}^2\le4H^2,
\]
even on a failed path.
The fresh-noise contribution is at most
$2\delta^2/(C_b\ell^2)$.  Thus every replicate has conditional second
moment at most $4H^2+\delta^2$, and
\eqref{eq:robust-global-mse} gives the bound $V_*$.

 An anchor costs $O(k\ell^2(N_\eps+P_\eps))$ calls and occurs at most
 $1+U/Q$ times.  A length-$h$ interval costs
 $O(k\ell^2(1+\min\{m,h^2,H^2/\delta^2\}+P_\eps))$ calls and is created
 at most $\lfloor U/h\rfloor$ times.  Indeed, $h$ divides the epoch length
 $Q$, and every creation of a length-$h$ object completes a new block of
 $h$ micro-edges.  Since the lengths are dyadic and
 $Q\asymp\sqrt{N_\eps}$,
\[
 \sum_{h\le Q,\ h\ {\rm dyadic}}\frac{h^2+1}{h}=O(Q),
 \qquad
 \sum_{h\le Q,\ h\ {\rm dyadic}}\frac1h=O(1).
\]
 The initial anchor is the only startup cost.  Altogether,
 \[
  S\le Ck\ell^2\left(N_\eps+P_\eps+
           U(\sqrt{N_\eps}+P_\eps)\right).
 \]
 Taking expectations proves \eqref{eq:dyadic-path-cost}.
\end{proof}

\subsection{Proof of Theorem~\ref{thm:dyadic-upper}}

\begin{proof}
Let ${\cal E}$ be the simultaneous-accuracy event from
Lemma~\ref{lem:dyadic-invariant}.  On ${\cal E}$, the direction perturbation
and descent inequalities give, pathwise,
\[
 f(x_{t+1})\le f(x_t)
 -\eta(1-\eta L_1/2)\norm{\nabla f(x_t)}
 +2\eta\delta+\eta^2L/2.
\]
Since $1-\eta L_1/2\ge7/8$, $\eta L/2\le\delta/8$, and
$\Delta/(\eta T)\le\delta/16$, summation yields
\begin{equation}
 \frac1T\sum_{t=0}^{T-1}\norm{\nabla f(x_t)}\le\frac52\delta,
 \qquad
 \sum_{t=0}^{T-1}\norm{\nabla f(x_t)}
       \le C(\Delta/\eta+T\delta).
 \label{eq:dyadic-average-gradient}
\end{equation}
The fresh uniform output therefore contributes at most $5\delta/2$ to
$\E[\norm{\nabla f(\widehat x)}\mathbf1_{\cal E}]$.

To control the complementary event, enumerate the constructed object errors
as $e_1,\ldots,e_{M_{\rm obj}}$, setting absent objects to zero.  The last
claim of Lemma~\ref{lem:dyadic-invariant} gives
\[
 \E\sum_{j=1}^{M_{\rm obj}}\norm{e_j}^2\le M_{\rm obj}V_*.
\]
Every accepted output point has $\norm{v_t}\le M$, and its estimator error
is the sum of at most $\ell$ object errors.  The abort output $x_0$ has
gradient norm at most $\mathcal G(\Delta)<M$.  Consequently, on every run,
\[
 \norm{\nabla f(\widehat x)}
 \le M+\left(\ell\sum_{j=1}^{M_{\rm obj}}\norm{e_j}^2\right)^{1/2}.
\]
Cauchy--Schwarz and \eqref{eq:dyadic-failure-budget} show that
\[
 \E[\norm{\nabla f(\widehat x)}\mathbf1_{{\cal E}^c}]
 \le Mp+\sqrt{p\ell M_{\rm obj}V_*}\le\delta/2.
\]
Thus the expected gradient norm is at most $3\delta<\eps$.

On ${\cal E}$, \eqref{eq:intrinsic-micro-path} and
\eqref{eq:dyadic-average-gradient} imply
\[
 U\le C\left(T+\frac{\eta L_1}{\delta}
                 \sum_t\norm{v_t}\right)
 \le C\left(T+\frac{L_1\Delta}{\delta}\right)
 \le C(1+A_\eps+B_\eps).
\]
On ${\cal E}^c$, use $U\le U_{\max}$.  Since $pU_{\max}\le1$,
\begin{equation}
 \E U\le C(1+A_\eps+B_\eps).
 \label{eq:virtual-path-length}
\end{equation}
Substitution into \eqref{eq:dyadic-path-cost} proves the sample bound.
The logarithmic exponent can be read directly from the parameters. Since
$\overline\Delta\le49\Delta+3\eta\delta$, one has
\begin{align*}
 M/\delta&\le C(1+\sqrt{A_\eps}+B_\eps),\\
 U_{\max}&\le C(1+A_\eps+B_\eps)^2,\\
 p^{-1}&\le C\ell^2(1+A_\eps+B_\eps)^2
                   (1+H^2/\eps^2).
\end{align*}
Thus $k,\ell\le C\mathcal L_\eps$ and
$k\ell^2\le C\mathcal L_\eps^3$. The cost estimate in
Lemma~\ref{lem:dyadic-invariant} therefore has precisely the logarithmic
upper factor used in Theorem~\ref{thm:main}.

All replicate means for the objects created at one virtual time can be
stacked and exactly averaged in a single polynomial-consensus operation.
Every agent can then apply the same finite robust aggregation rule locally.
Thus $O(U+1)$ exact averaging operations suffice.  Lemma
\ref{lem:polynomial-consensus} and \eqref{eq:virtual-path-length} give
\eqref{eq:dyadic-communication}.
\end{proof}

The reference-point organization is related to nested variance reduction
\cite{snvrg2018}.  The present construction separately averages fresh oracle
noise at each endpoint, and its analysis uses simultaneous robust accuracy
instead of same-sample smoothness or unbiasedness at adaptively selected
virtual times.

\section{A retiring chain with uniformly bounded disagreement}
\label{subsec:retiring-dwh}

Dense weak hiding spreads a hidden direction across component rows, so that
one row is weakly informative while the average preserves the signal.
We use the sign-table and radial-coordinate construction of
Peng et al.~\cite{dense2026}. The additional requirement here is a uniform
bound on local gradient disagreement, independent of the number of stages.
A retiring selector removes the component-specific signal of a completed
stage and passes it to the next one. The following calculus lemma controls
both its gradient and its Hessian without dependence on the chain length.

\subsection{A gate with a square-root plateau bound}

Keep the fixed thresholds
\[
 a=\frac5{32},\qquad b=\frac9{32}.
\]
For $u=(s-a)/(b-a)$ define
\begin{equation}
 G(s)=
 \begin{cases}
 0,&s\le a,\\
 10u^3-15u^4+6u^5,&a<s<b,\\
 1,&s\ge b.
 \end{cases}
 \label{eq:retiring-gate}
\end{equation}
This is the standard quintic smoothstep.

\begin{lemma}[Stable plateau gate]
\label{lem:stable-gate}
The function in \eqref{eq:retiring-gate} is nondecreasing and $C^2$.  There
are universal constants $M_1,M_2,\kappa_G$ such that
\begin{equation}
 0\le G\le1,\qquad
 \norm{G'}_\infty\le M_1,\qquad
 \norm{G''}_\infty\le M_2,\qquad
 |G'(s)|^2\le\kappa_G(1-G(s)).
 \label{eq:gate-square-root}
\end{equation}
Moreover, $G'=G''=0$ on both plateaus.
\end{lemma}

\begin{proof}
Only the last inequality needs comment.  On $u\in(0,1)$,
\[
 G'(s)=\frac{30}{b-a}u^2(1-u)^2,
 \qquad
 1-G(s)=(1-u)^3(1+3u+6u^2).
\]
Hence
\[
 \frac{|G'(s)|^2}{1-G(s)}
 =\frac{900}{(b-a)^2}
   \frac{u^4(1-u)}{1+3u+6u^2}
 \le \frac{900}{(b-a)^2}.
\]
The endpoint derivatives of orders one and two vanish, so the polynomial
piece joins the constant plateaus in $C^2$.
\end{proof}

\subsection{A dimension-free stick-breaking calculus lemma}

The next lemma is the mechanism that prevents accumulation over the chain
length.  The variables $x_j$ below live in mutually orthogonal Euclidean
blocks.

\begin{lemma}[Uniform $C^2$ stability of a retiring selector]
\label{lem:stick-breaking}
For $j=1,\ldots,T$, let $d_j=d_j(x_j)$ and
$g_j=g_j(x_j)\in[0,1]$, and recursively define
\begin{equation}
 Z_{T+1}=0,
 \qquad
 Z_j=(1-g_j)d_j+g_jZ_{j+1}.
 \label{eq:stick-recursion}
\end{equation}
Assume, uniformly in $j$,
\begin{align}
 |d_j|&\le D_0,&
 \norm{\nabla d_j}&\le D_1,&
 \norm{\nabla^2d_j}_{\rm op}&\le D_2,
 \label{eq:stick-d-bounds}\\
 \norm{\nabla g_j}^2&\le\kappa(1-g_j),&
 \norm{\nabla^2g_j}_{\rm op}&\le M_2,
 \label{eq:stick-g-bounds}
\end{align}
and suppose $g_j=1$ implies
$\nabla g_j=\nabla^2g_j=0$.  Then there are constants
$C_1,C_2$, depending only on
$(D_0,D_1,D_2,\kappa,M_2)$ and not on $T$, such that
\begin{equation}
 |Z_j|\le D_0,
 \qquad
 \norm{\nabla Z_j}\le C_1,
 \qquad
 \norm{\nabla^2Z_j}_{\rm op}\le C_2
 \quad(1\le j\le T+1).
 \label{eq:stick-uniform}
\end{equation}
More precisely, if $D_k=O(K)$ for $k=0,1,2$ and
$\kappa,M_2=O(1)$, then $C_1,C_2=O(K)$.
\end{lemma}

\begin{proof}
The value bound follows by backward induction because
\eqref{eq:stick-recursion} is a convex combination of $d_j$ and $Z_{j+1}$.
Write $g=g_j$, $d=d_j$, and $V=Z_{j+1}$.  The current-block and tail parts of
the gradient are orthogonal and equal to
\[
 a=(1-g)\nabla d+(V-d)\nabla g,
 \qquad
 g\nabla V.
\]
Consequently,
\[
 \norm a^2
 \le \bigl(2D_1^2+8\kappa D_0^2\bigr)(1-g)
 =:K_1^2(1-g).
\]
If $\norm{\nabla V}\le K_1$, then
\[
 \norm{\nabla Z_j}^2
 \le K_1^2(1-g)+g^2K_1^2
 \le K_1^2,
\]
which closes the gradient induction.

For the Hessian, with respect to the decomposition
$x_j\oplus(x_{j+1},\ldots,x_T)$,
\begin{equation}
 \nabla^2Z_j=
 \begin{pmatrix}
 A&B\\ B^\top&g\nabla^2V
 \end{pmatrix},
 \label{eq:stick-hessian-block}
\end{equation}
where
\begin{align*}
 A&=(1-g)\nabla^2d
 -\nabla g\otimes\nabla d-\nabla d\otimes\nabla g
 +(V-d)\nabla^2g,\\
 B&=\nabla g\otimes\nabla V.
\end{align*}
Thus
\[
 \norm{A}_{\rm op}\le
 D_2+2\sqrt\kappa D_1+2D_0M_2=:A_*,
 \qquad
 \norm{B}_{\rm op}^2\le\kappa(1-g)K_1^2.
\]
Choose
\[
 K_2\ge 2A_*+2\sqrt\kappa K_1.
\]
Assume inductively that $\norm{\nabla^2V}_{\rm op}\le K_2$.  For $g<1$,
the lower-right block of either matrix
$K_2I\mp\nabla^2Z_j$ is bounded below by
$K_2(1-g)I$.  Its Schur complement is therefore bounded below by
\[
 \left(K_2-A_*-\frac{\kappa K_1^2}{K_2}\right)I\succeq0.
\]
If $g=1$, then $A=B=0$ by the plateau assumption and the conclusion follows
directly from the induction hypothesis.  Hence
$-K_2I\preceq\nabla^2Z_j\preceq K_2I$, which proves the Hessian bound.  The
last assertion follows from the displayed choices of $K_1,K_2$.
\end{proof}

\subsection{The retiring radial chain}

We use the same dense weak-hiding table as in \cite{dense2026}.  For $m$
agents its admissible bias grid is
\begin{equation}
 \mathcal R_m:=\left\{
 \frac{k}{m}: k\in\mathbb Z_{>0},\ k\equiv m\pmod 2,\
 k\le\frac m8
 \right\}.
 \label{eq:retiring-bias-grid}
\end{equation}
For a bias $\rho\in\mathcal R_m$, stage $j$ has an independent hidden
direction $\Theta^{(j)}\in\{\pm1\}^D$.  Conditional on
$\Theta^{(j)}$, the columns of the table
$S^{(j)}\in\{\pm1\}^{m\times D}$ are independent.  In column $h$, a
uniformly random subset of $m(1+\rho)/2$ rows receives the sign
$\Theta_h^{(j)}$ and all remaining rows receive the opposite sign.  Thus the
row average is exactly $\rho\Theta^{(j)}$.  Consecutive admissible biases
differ by $2/m$.  Let
\begin{equation}
 \Psi(z)=\frac{z}{\sqrt{1+\norm z^2}},
 \quad
 q_j(z)=\frac1{\sqrt D}\inner{\Theta^{(j)}}{\Psi(z)},
 \quad
 \widetilde q_{ij}(z)=
 \frac1\rho\inner{S^{(j)}_{i:}/\sqrt D}{\Psi(z)}.
 \label{eq:retiring-radial-coordinates}
\end{equation}
Then $m^{-1}\sum_i\widetilde q_{ij}=q_j$ and, uniformly in $i,j,z$,
\begin{equation}
 |q_j|\le1,
 \quad \norm{\nabla q_j}\le1,
 \quad \norm{\nabla^2q_j}_{\rm op}\le2,
 \qquad
 |\widetilde q_{ij}|\le\rho^{-1},
 \quad \norm{\nabla\widetilde q_{ij}}\le\rho^{-1},
 \quad \norm{\nabla^2\widetilde q_{ij}}_{\rm op}\le2\rho^{-1}.
 \label{eq:retiring-radial-bounds}
\end{equation}
Put $g_j=G(q_j(x_j))$ and define the component selectors
\begin{equation}
 Z_{i,T+1}=0,
 \qquad
 Z_{i,j}=(1-g_j)\widetilde q_{ij}(x_j)+g_jZ_{i,j+1}.
 \label{eq:component-retiring-selector}
\end{equation}
Their average $Z_j=m^{-1}\sum_iZ_{i,j}$ obeys the same recursion with
$\widetilde q_{ij}$ replaced by $q_j$.  Fix
\[
 C=2,
 \qquad
 \lambda=\frac{33}{16},
\]
and define
\begin{equation}
 R_i(x)=
 -Z_{i,1}(x)
 -\sum_{j=1}^Tg_j(x_j)\bigl(q_j(x_j)+C\bigr)
 +\frac\lambda2\sum_{j=1}^T\norm{x_j}^2,
 \qquad
 R=\frac1m\sum_{i=1}^mR_i.
 \label{eq:retiring-chain}
\end{equation}

\begin{lemma}[Tunable-bias sequential revelation]
\label{lem:tunable-sequential-revelation}
For the independent exact-bias stage distribution above, let
\begin{equation}
 r_\rho=
 \left\lfloor\frac1{4096}\min\{m,\rho^{-2}\}\right\rfloor.
 \label{eq:retiring-row-threshold}
\end{equation}
Consider an augmented interaction with the following properties.
\begin{enumerate}
\item The data available immediately before a stage's initialization are
independent of its stage pair.  Initialization may reveal one table row,
whose index is selected before that row is seen.
\item Before completion, each full value--gradient reply is a measurable
function of the past, the query, completed stage pairs, and the selected
rows of the current stage.  This formula does not itself test the hidden
alignment and contains no future-stage data.
\item Completion is checked before the reply, at the first query whose
current-stage alignment exceeds $a$.  A completion response may reveal the
completed stage pair and one initialization row of the successor.  Each
call is assigned only to the stage current before that call.
\item There is a prefix-truncated interaction satisfying the preceding
properties.  It has the same replies as the full interaction whenever all
stages beyond the current one have alignment at most $a$.  This equality
includes the final output.
\end{enumerate}
Let $1\le J<T$, $M\ge1$, and $r_\rho\ge2$.  If
\begin{equation}
 M\le\frac{r_\rho J}{2},\qquad
 D\ge\max\left\{
 512\log(M+3)+2048,
 \frac{2048}{25}\log\bigl(128(M+1)T\bigr)
 \right\},
 \label{eq:imported-sequential-budget}
\end{equation}
then
\begin{equation}
 \Pr(\nu_M<J,\mathcal E_{\rm future}^{c})\ge\frac{95}{128}.
 \label{eq:imported-sequential-probability}
\end{equation}
Here $\nu_M$ counts completed stages in the truncated interaction, and
$\mathcal E_{\rm future}$ is the event that some query or the final output
has alignment greater than $a$ at a stage beyond its current stage.
On the displayed event, the full-oracle output satisfies
$q_k(\widehat x_k)\le a$ for every $k>J$.
\end{lemma}

\begin{proof}
We give the information argument explicitly.  It is the dense-weak-hiding
principle of~\cite{dense2026}, with the bias retained as a free parameter.

\emph{One stage.}
Suppress the stage index.  Condition on any side information $W$ independent
of $(\Theta,S)$, including the private random seed and the data available
before initialization.  Grant the algorithm the entire selected row on
every call.  This can only strengthen it.  Consider a row-only continuation
with at most $B=M+1$ predictable query vectors $V_t$, $\|V_t\|\le1$.
After the $r_\rho$th distinct row has been returned, replace subsequent
query vectors by zero and stop returning new rows.  Let $\mathcal Y$ be
this stopped transcript, including the optional initialization row and all
query vectors.  It reads at most $r_\rho$ distinct rows.

Suppose $d<r_\rho$ distinct rows have been seen, and write
$A_h$ for their sum in column $h$.  Conditional exchangeability of the
unseen row labels gives
\[
 \Pr(S_{ih}=1\mid\Theta,\text{past})
 =\frac{m+m\rho\Theta_h-d-A_h}{2(m-d)}=:p_h(\Theta_h)
\]
for each newly selected row.  The same formula holds after earlier
coordinates of that new row are disclosed, because table columns are
independent conditional on $\Theta$.  Since $\rho\le1/8$ and
$d\le m/4096$, both possible values of $p_h$ lie in $[1/4,3/4]$ and differ
by $m\rho/(m-d)$.  The binary entropy function has second derivative of
absolute value at most $16/3$ on this interval.  Its Jensen gap therefore
bounds the information in one coordinate by
\[
 I(\Theta;S_{ih}\mid\text{past},S_{i,<h})
 \le \frac23\left(\frac{m\rho}{m-d}\right)^2
 \le\rho^2.
\]
Repeated rows contain no additional information, and query choices are
predictable from the preceding transcript.  The chain rule yields
\[
 I(\Theta;\mathcal Y\mid W)
 \le D r_\rho\rho^2\le\frac{D}{4096}.
\]

Let $E$ be the event that a predictable query made before the row cutoff
has $\langle\Theta,V_t\rangle/\sqrt D>a$.
Under the product law that keeps the transcript distribution but makes
$\Theta$ independent of it, a Rademacher tail bound gives
$\Pr(E)\le B\exp(-a^2D/2)$.  Binary data processing for relative entropy
therefore gives, under the true joint law,
\[
 \Pr(E\mid W)
 \le\frac{D/4096+\log2}{a^2D/2-\log B}
 \le\frac18.
\]
For the last inequality use $a=5/32$,
$B\le M+1$, and the first dimension condition
in~\eqref{eq:imported-sequential-budget}.
All these estimates are uniform in the independent side information.

To apply this estimate to a stage of the actual truncated interaction,
continue its precompletion reply formula without ever announcing
completion.  The resulting row-only interaction has the same histories
and queries through its first completion query.  Completion is determined
before that query's response, so completion in fewer than $r_\rho$
assigned calls uses fewer than $r_\rho$ previously seen rows, even with
one initialization row.  It is consequently contained in $E$ and has
conditional probability at most $1/8$.

\emph{Future stages.}
The prefix-truncated interaction reads only completed pairs and the current
stage pair.  It is defined on every realization, without conditioning on
the absence of a bypass.  Every direction beyond its current stage is thus
independent of its query block.  The same holds at the final output.
A Rademacher tail bound and the second dimension condition give
\[
 \Pr(\mathcal E_{\rm future})
 \le T(M+1)\exp(-a^2D/2)\le\frac1{128}.
\]
The transcript-matching assumption couples the full and truncated
interactions on the complementary event.

\emph{Counting calls.}
Let $Z$ count the stages among the first $J$ that are reached and complete
in fewer than $r_\rho$ assigned calls.  Conditioning immediately before
each reached stage is initialized and using the one-stage bound gives
$\E Z\le J/8$.  If all first $J$ stages complete within
$M\le r_\rho J/2$ calls, then $Z\ge J/2$, because assigned calls for
different stages are disjoint.  Markov's inequality yields
$\Pr(\nu_M\ge J)\le1/4$.  Combining this with the future-stage bound
gives $1-1/4-1/128=95/128$.  Finally, on $\{\nu_M<J\}$ every $k>J$
is beyond the current stage; the output part of
$\mathcal E_{\rm future}^{c}$ gives the claimed final alignments.
\end{proof}


\begin{proposition}[Analytic and oracle properties of the retiring chain]
\label{prop:retiring-chain-properties}
There are universal constants $C_{\rm sm},C_{\rm dis},C_{\rm gap},\chi>0$
such that, for every $T$, $m$, admissible $\rho$, and every realization of
the hidden tables, the chain \eqref{eq:retiring-chain} has the following
properties.
\begin{enumerate}
\item Each component is ordinarily smooth with
\begin{equation}
 \sup_x\norm{\nabla^2R_i(x)}_{\rm op}\le\frac{C_{\rm sm}}\rho.
 \label{eq:retiring-smoothness}
\end{equation}
\item The component--average disagreement is uniform in both $x$ and $T$:
\begin{equation}
 \sup_x\norm{\nabla R_i(x)-\nabla R(x)}
 \le\frac{C_{\rm dis}}\rho.
 \label{eq:retiring-disagreement}
\end{equation}
\item The initial gap satisfies
\begin{equation}
 R(0)-\inf R\le C_{\rm gap}T.
 \label{eq:retiring-gap}
\end{equation}
\item If at least one stage is not fully open, then
\begin{equation}
 \min_{1\le j\le T}q_j(x_j)<b
 \quad\Longrightarrow\quad
 \norm{\nabla R(x)}\ge\chi.
 \label{eq:retiring-gradient-lower}
\end{equation}
\item The exact component value--gradient oracle satisfies the
interface conditions of Lemma~\ref{lem:tunable-sequential-revelation}.
\end{enumerate}
\end{proposition}

\begin{proof}
For item 1, Lemma~\ref{lem:stable-gate} and
\eqref{eq:retiring-radial-bounds} give
\[
 \norm{\nabla g_j}^2\le\kappa_G(1-g_j),
 \qquad
 \norm{\nabla^2g_j}_{\rm op}\le C_G
\]
with universal $C_G$.  Apply Lemma~\ref{lem:stick-breaking} to
\eqref{eq:component-retiring-selector} with
$(D_0,D_1,D_2)=\rho^{-1}(1,1,2)$.  It yields
\[
 \sup_x\norm{\nabla^2Z_{i,1}(x)}_{\rm op}=O(\rho^{-1}).
\]
Each common term $g_j(q_j+C)$ depends on only one block and has a uniformly
bounded Hessian.  Their Hessians form a block-diagonal matrix, so their sum
has operator norm $O(1)$ rather than $O(T)$.  Adding the quadratic proves
\eqref{eq:retiring-smoothness}.

For item 2, let $D_{i,j}=Z_{i,j}-Z_j$.  Then
\begin{equation}
 D_{i,T+1}=0,
 \qquad
 D_{i,j}=(1-g_j)(\widetilde q_{ij}-q_j)+g_jD_{i,j+1}.
 \label{eq:retiring-difference-recursion}
\end{equation}
The three $C^2$ bounds on $\widetilde q_{ij}-q_j$ are all $O(\rho^{-1})$.
Lemma~\ref{lem:stick-breaking} therefore gives
$\sup_x\norm{\nabla D_{i,1}(x)}=O(\rho^{-1})$.  Since all other terms in
\eqref{eq:retiring-chain} are common to the components,
$R_i-R=-D_{i,1}$, proving \eqref{eq:retiring-disagreement}.

At the origin, all $q_j$ and all gates vanish, so $R(0)=0$.  Moreover,
$|Z_1|\le1$, $0\le g_j\le1$, $|q_j|\le1$, and the quadratic is nonnegative.
Thus
\[
 R(x)\ge-1-(C+1)T,
\]
which proves item 3.

For item 4, let $j$ be the first index with $q_j(x_j)<b$.  All preceding
gates equal one, hence $Z_1=Z_j$.  The recursion gives the exact identity
\begin{equation}
 Z_j+g_j(q_j+C)=q_j+g_j(C+Z_{j+1}).
 \label{eq:retiring-gradient-identity}
\end{equation}
Because $|Z_{j+1}|\le1$, $C=2$, and $G'\ge0$, differentiation in block $j$
gives
\[
 \nabla_{x_j}R(x)=\lambda x_j-c_j\nabla q_j(x_j),
 \qquad
 c_j=1+G'(q_j)(C+Z_{j+1})\ge1.
\]
The radial estimate used in Lemma~A.6 of \cite{dense2026}, with the same
$\lambda$ and threshold $b$, gives
$\norm{\lambda x_j-c_j\nabla q_j(x_j)}\ge\chi$ uniformly for
$c_j\ge1$ and $q_j<b$.  This proves \eqref{eq:retiring-gradient-lower}.

For item 5, define the prefix-truncated interaction explicitly.
Reveal each completed stage pair in full, which can only strengthen the
algorithm. If $h$ stages have been completed, stage $k=h+1$ is current.
On a query to component $i$, the truncated oracle retains the completed
stage data and the public quadratic term. It evaluates the current-stage
completion test before returning the reply.

If $q_k\le a$, set every gate with index at least $k$ to zero in the
simulated value and gradient, and terminate the selector with
$Z_{i,k}=\widetilde q_{ik}$. The resulting reply depends only on the
completed stage pairs, the query, and the selected current row
$S^{(k)}_{i:}$. Its algebraic precompletion formula contains no test of
$q_k$ and no unread future-stage data; the completion test is used only to
choose when to leave this formula.

If $q_k>a$, disclose the completed stage pair, retain its actual gate, and
initialize the successor by the selected row $S^{(k+1)}_{i:}$ when $k<T$.
In this reply, force every gate with index at least $k+1$ to zero and set
$Z_{i,k+1}=\widetilde q_{i,k+1}$. Thus the completion reply uses the newly
completed pair and at most one successor row. The physical call remains
assigned to stage $k$; the successor row is the optional initialization
row allowed by Lemma~\ref{lem:tunable-sequential-revelation}.

This construction defines the simulated interaction on every realization.
It reads only completed stage pairs and the current pair, with the
successor row first read on its initialization. Future pairs are independent
of the history immediately before that initialization. On the event that
all stages beyond the current one satisfy $q_j\le a$, their gates and first
derivatives vanish. Consequently all omitted terms in the common reward
sum are zero, and the full selector cuts off at the same current or
successor block as the simulated selector. The simulated and full component
values and every gradient block are then identical. Induction gives
pathwise agreement of the histories, queries, and final outputs on the
no-future-alignment event. These are precisely the four properties in
Lemma~\ref{lem:tunable-sequential-revelation}.
\end{proof}

\subsection{Proof of Theorem~\ref{thm:retiring-dwh-lower}}

\begin{proof}
Set
\[
 L_\star=L_0+\frac12L_1H,
 \qquad
 A_\star=\frac{L_\star\Delta}{\eps^2},
 \qquad
 K:=\min\left\{m,\frac{H^2}{\eps^2}\right\}.
\]
Then $L_H/2\le L_\star\le L_H$, so $A_\star=\Theta(A_\eps)$ and
$\sqrt{N_\eps}=\Theta(1+\sqrt K)$.  We first treat the nonboundary case
\begin{equation}
 m\ge m_0,\qquad K\ge K_0,\qquad
 A_\eps\ge C_0\sqrt K,
 \label{eq:retiring-interior-case}
\end{equation}
for sufficiently large universal constants.  The complementary cases will be
absorbed by the earlier centralized and discovery lower bounds.

For an admissible bias $\rho$ scale the retiring chain by
\begin{equation}
 \beta_\rho=\frac{L_\star\rho\chi}
 {8C_{\rm sm}\eps},
 \qquad
 \alpha_\rho=\frac{64C_{\rm sm}\eps^2}
 {L_\star\rho\chi^2},
 \qquad
 f_i(y)=\alpha_\rho R_i(\beta_\rho y).
 \label{eq:retiring-scaling}
\end{equation}
Then
\[
 \alpha_\rho\beta_\rho^2\frac{C_{\rm sm}}\rho=L_\star,
 \qquad
 \alpha_\rho\beta_\rho\chi=8\eps.
\]
Taking
\begin{equation}
 T=\left\lfloor c_TA_\star\rho\right\rfloor
 \label{eq:retiring-stage-count}
\end{equation}
with a sufficiently small universal $c_T$ makes the initial gap at most
$\Delta$, while Proposition~\ref{prop:retiring-chain-properties} gives
ordinary $L_\star$-smoothness and an $8\eps$ gradient lower bound until all
stages are fully open.

The scaled uniform disagreement is
\begin{equation}
 \sup_y\norm{\nabla f_i(y)-\nabla f(y)}
 \le\alpha_\rho\beta_\rho\frac{C_{\rm dis}}\rho
 =\frac{8C_{\rm dis}}\chi\frac{\eps}{\rho}.
 \label{eq:retiring-scaled-disagreement}
\end{equation}
Let
\begin{equation}
 \bar\rho=C_\rho\max\left\{m^{-1/2},\frac\eps H\right\},
 \label{eq:retiring-bias-choice}
\end{equation}
where $C_\rho$ is a sufficiently large universal constant.  In the interior
case we may assume $\bar\rho\le1/16$: otherwise either $m$ or $H/\eps$ is
bounded by a universal constant, hence $K=O(1)$, one of the boundary cases
below.  Choose the least $\rho\in\mathcal R_m$ with $\rho\ge\bar\rho$.
Since consecutive values of $\mathcal R_m$ differ by $2/m$ and
$\bar\rho\ge C_\rho/\sqrt m$, increasing $m_0$ if necessary gives
\begin{equation}
 \bar\rho\le\rho\le\bar\rho+\frac2m\le2\bar\rho\le\frac18.
 \label{eq:retiring-grid-rounding}
\end{equation}
Thus \eqref{eq:retiring-scaled-disagreement} is at most $H/2$, after fixing
$C_\rho$ large enough.

Append orthogonal simplex vectors $u_i$ of norm $H/2$ and mean zero, and set
\begin{equation}
 F_i(y,t)=f_i(y)+\inner{u_i}{t}.
 \label{eq:retiring-relaxed-embedding}
\end{equation}
The total local--global disagreement is at most $H/\sqrt2<H$.  Every $F_i$
is ordinarily $L_\star$-smooth and
$\norm{\nabla F_i}\ge H/2$, so
\[
 L_\star=L_0+\frac12L_1H
 \le L_0+L_1\norm{\nabla F_i}.
\]
Hence the required local relaxed smoothness holds globally.  The auxiliary
linear terms average to zero, so the average objective, its initial gap, and
its stationarity obstruction are unchanged.  The exact-gradient oracle used
here has variance zero and is therefore an admissible subclass for every
stated value of $\sigma$.

For the bias in \eqref{eq:retiring-grid-rounding}, a two-case calculation gives
\begin{equation}
 \min\{m\rho,\rho^{-1}\}
 =\Theta(\sqrt K).
 \label{eq:retiring-effective-size}
\end{equation}
Indeed, when $H/\eps\le\sqrt m$ one has
$\rho=\Theta(\eps/H)$ and the minimum is $\Theta(H/\eps)$; otherwise
$\rho=\Theta(m^{-1/2})$ and both branches are $\Theta(\sqrt m)$.
The tunable-bias one-stage row threshold is
$r_\rho=\Theta(\min\{m,\rho^{-2}\})$, while
$T=\Theta(A_\star\rho)$.  Conditions
\eqref{eq:retiring-interior-case} ensure $r_\rho,T\ge2$, and hence their
product is
\begin{equation}
 r_\rho T
 =\Theta\left(
 A_\star\min\{m\rho,\rho^{-1}\}
 \right)
 =\Theta(A_\eps\sqrt K).
 \label{eq:retiring-row-stage-product}
\end{equation}
For any fixed call budget $M$ below a sufficiently small universal multiple
of this quantity, choose a finite block dimension
\begin{equation}
 D\ge C_D\left[1+\log\bigl((M+3)(T+1)\bigr)\right].
 \label{eq:retiring-dimension}
\end{equation}
Take $J=T-1$ and restrict $M$ so that
$M\le r_\rho(T-1)/2$.  Proposition~\ref{prop:retiring-chain-properties}
verifies the sequential interface, so Lemma~\ref{lem:tunable-sequential-revelation} gives
an event of probability at least $95/128$ on which
$q_T(\widehat x_T)\le a<b$.  By the deterministic $8\eps$ stationarity
obstruction, conditioning on a randomized algorithm's private seed and then
averaging over the stage pairs fixes one realization for which
\begin{equation}
 \Pr\!\left(\norm{\nabla F(\widehat x)}\ge8\eps\right)
 \ge\frac{95}{128}.
 \label{eq:retiring-fixed-budget-failure}
\end{equation}

We now transfer the statement both to decentralized protocols and to expected
call budgets.  A centralized IFO simulator maintains the complete state of all
agents, executes every local arithmetic and communication operation for free,
and replaces each local exact-gradient query by one query to the corresponding
component.  Simultaneous batches can be serialized without changing their
responses.  Thus a decentralized protocol making $S$ total local calls induces
an IFO algorithm making exactly $S$ calls; we apply the lower bound to the
output of any fixed agent.  The information argument above allows an oracle returning both
value and gradient. Its lower bound therefore applies to the weaker
gradient-only protocol considered here.

To make the dimension quantifier explicit, choose first
\[
 M_0=\left\lfloor c_Mr_\rho(T-1)\right\rfloor
\]
with a sufficiently small universal $c_M>0$, and choose $D$ using
\eqref{eq:retiring-dimension} with $M=M_0$. These choices depend only on the
problem parameters. Increasing the universal interior thresholds if needed
ensures $M_0\ge1$ and $M_0=\Theta(A_\eps\sqrt K)$.
Consider any protocol in this fixed dimension whose worst-case expected
call count is $q\le M_0/4$. Cap its simulator after $M_0$ calls, using a
fixed fallback output if necessary. The fixed-budget argument supplies one
hard instance for which the capped output has gradient at least $8\eps$
with probability at least $95/128$. On that same instance, Markov's
inequality gives $\Pr(S>M_0)\le1/4$. Therefore the original output satisfies
\[
 \Pr\!\left(\norm{\nabla F(\widehat x)}\ge8\eps\right)
 \ge\frac{95}{128}-\frac14=\frac{63}{128},
\]
and its expected gradient norm exceeds $\eps$. This contradiction proves
$q=\Omega(A_\eps\sqrt K)$ with the dimension fixed before the algorithm
is selected.

Finally consider the omitted cases.  If $A_\eps\lesssim\sqrt{N_\eps}$, the
hidden-sign lower bound $\Omega(N_\eps)$ dominates
$A_\eps\sqrt{N_\eps}$.  If $K$, $m$, or $H/\eps$ is bounded by a universal
constant, the ordinary centralized lower bound $\Omega(A_\eps)$ dominates the
same target.  These observations also cover $H=0$, an empty admissible grid,
$r_\rho<2$, $T<2$, and $\bar\rho>1/16$.  Since
$\sqrt{N_\eps}=\Theta(1+\sqrt K)$, combining the interior and boundary cases
proves \eqref{eq:retiring-square-root-lower}.
\end{proof}

\section{A terminally flat chain for the coarse stochastic lower bound}
\label{sec:stochastic-coarse-closure}

The construction uses a bounded hidden gradient in each block and a public
radial wall. Terminal flatness prevents residual curvature from violating
the relaxed-smoothness condition near a stationary point. A Bernoulli oracle
then reveals at most one new hidden direction per success. The rare-reveal
principle is related to~\cite{arjevani2022}; the analytic certificate and
coupling below verify its use in the present function class.

\subsection{A flat switch with approximate logarithmic derivatives}

Define
\begin{equation}
 \vartheta(t)=
 \begin{cases}
  0,&t\le0,\\
  e^{-1/t},&t>0,
 \end{cases}
 \qquad
 \mathsf S(t)=
 \frac{\vartheta(t)}{\vartheta(t)+\vartheta(1-t)}.
 \label{eq:flat-switch}
\end{equation}
Then $\mathsf S$ is nondecreasing and $C^\infty$, equals zero on
$(-\infty,0]$, equals one on $[1,\infty)$, and satisfies
$\mathsf S(1-t)=1-\mathsf S(t)$.

\begin{lemma}[Flat-switch calculus]
\label{lem:flat-switch-calculus}
There is a universal constant $C_{\rm sw}$ such that, for every
$0<\zeta\le e^{-4}$, with
\begin{equation}
 \Gamma_\zeta
 :=C_{\rm sw}\bigl(1+\log(1/\zeta)\bigr)^4,
 \label{eq:flat-switch-Gamma}
\end{equation}
one has, for every $t\in\R$,
\begin{align}
 0\le \mathsf S'(t)
 &\le \Gamma_\zeta\bigl(\zeta+\mathsf S(t)\bigr),
 \label{eq:switch-first-left}\\
 \mathsf S'(t)
 &\le \Gamma_\zeta\bigl(\zeta+1-\mathsf S(t)\bigr),
 \label{eq:switch-first-right}\\
 |\mathsf S''(t)|
 &\le \Gamma_\zeta\bigl(\zeta+\mathsf S'(t)\bigr).
 \label{eq:switch-second}
\end{align}
The same statements hold after any fixed affine rescaling of the argument,
with only a universal change in $C_{\rm sw}$.
\end{lemma}

\begin{proof}
Outside $(0,1)$ all derivatives vanish.  On $(0,1/2]$, put
\[
 A(t)=\frac1t-\frac1{1-t},
 \qquad
 B(t)=\frac1{t^2}+\frac1{(1-t)^2}.
\]
Then
\[
 \mathsf S(t)=\frac1{1+e^{A(t)}},
 \qquad
 \mathsf S'(t)=\mathsf S(t)(1-\mathsf S(t))B(t),
\]
and direct differentiation gives
\begin{equation}
 |\mathsf S''(t)|
 \le C\mathsf S(t)(t^{-4}+1).
 \label{eq:flat-switch-second-raw}
\end{equation}
Moreover, on this half-interval
\begin{equation}
 A(t)\ge t^{-1}-2,
 \qquad
 \frac12e^{-A(t)}\le\mathsf S(t)\le e^{-A(t)},
 \qquad
 \mathsf S'(t)\ge4\mathsf S(t).
 \label{eq:flat-switch-elementary}
\end{equation}

Write $L=1+\log(1/\zeta)$.  If $\mathsf S(t)\ge\zeta$, then
$A(t)\le\log(1/\zeta)$, so \eqref{eq:flat-switch-elementary} yields
$t^{-1}\le C L$.  Hence
\[
 \mathsf S'(t)\le C L^2\mathsf S(t),
 \qquad
 |\mathsf S''(t)|\le C L^4\mathsf S(t)
 \le C L^4\mathsf S'(t).
\]
If $\mathsf S(t)<\zeta$, then by
\eqref{eq:flat-switch-elementary}, with $A=A(t)$,
\[
 \mathsf S(t)t^{-k}
 \le e^{-A}(A+2)^k
 \le C\zeta L^k,
 \qquad k=2,4.
\]
The last inequality is the elementary monotonicity bound for
$e^{-A}(A+2)^k$ on
$A\ge\log(1/(2\zeta))$, using $\zeta\le e^{-4}$ and enlarging a universal
constant.  Together with \eqref{eq:flat-switch-second-raw}, these two cases
prove \eqref{eq:switch-first-left} and \eqref{eq:switch-second} on
$(0,1/2]$.

On $[1/2,1)$, the first-left inequality is immediate from boundedness of
$\mathsf S'$ and $\mathsf S\ge1/2$.  The identity
$\mathsf S(1-t)=1-\mathsf S(t)$ transfers the first-left estimate on
$(0,1/2]$ to \eqref{eq:switch-first-right} on $[1/2,1)$ and transfers the
second-derivative estimate to the other half.  The remaining halves of the
two first-derivative estimates are again trivial by boundedness.  Fixed
affine rescalings multiply the first and second derivatives by fixed
constants, which are absorbed into $C_{\rm sw}$.
\end{proof}

Fix once and for all
\begin{equation}
 a=\frac1{16},
 \qquad b=\frac18,
 \qquad d=\frac14.
 \label{eq:coarse-chain-thresholds}
\end{equation}
Define the activation gate
\begin{equation}
 s(u)=\mathsf S\!\left(\frac{u-a}{b-a}\right)
 \label{eq:coarse-chain-gate}
\end{equation}
and the bounded reward
\begin{equation}
 r(u)=\int_{-\infty}^{u}
 \mathsf S(v+2)
 \left[1-\mathsf S\!\left(\frac{v-b}{d-b}\right)\right]dv.
 \label{eq:coarse-chain-reward}
\end{equation}
Write $R=\sup_u r(u)<\infty$.  On the interval $(-1,1)$,
\begin{equation}
 r'(u)=1\quad(u\le b),
 \qquad
 r'(u)=0\quad(u\ge d),
 \qquad
 0\le r'(u)\le1.
 \label{eq:reward-plateaux}
\end{equation}
Moreover, Lemma~\ref{lem:flat-switch-calculus} gives
\begin{equation}
 |r''(u)|\le C\Gamma_\zeta(\zeta+r'(u)),
 \qquad u\in(-1,1).
 \label{eq:reward-log-derivative}
\end{equation}

\subsection{A terminally flat scalar zero-chain}

For $q=(q_1,\ldots,q_T)\in(-1,1)^T$, define
\begin{equation}
 \Phi_T(q)
 =-r(q_1)-\sum_{j=2}^{T}s(q_{j-1})r(q_j).
 \label{eq:terminal-flat-scalar-chain}
\end{equation}
Use the conventions $s(q_0)=1$ and $r(q_{T+1})=0$.  Its coordinate
gradients have one sign:
\begin{equation}
 -\partial_j\Phi_T(q)
 =s(q_{j-1})r'(q_j)+s'(q_j)r(q_{j+1})
 =:g_j(q)\ge0.
 \label{eq:scalar-chain-gradient}
\end{equation}
Thus the two mechanisms in one coordinate never cancel.

\begin{lemma}[Scalar-chain properties]
\label{lem:scalar-chain-properties}
There are universal constants $G_0,C_0>0$, independent of $T$, such that:
\begin{enumerate}
\item $0\le g_j(q)\le G_0$ for every $j,q$;
\item if $q_j\le b$ for at least one $j$, then
      $\norm{\nabla\Phi_T(q)}\ge1$;
\item
\begin{equation}
 \norm{\nabla^2\Phi_T(q)}_{\rm op}
 \le C_0\Gamma_\zeta
       \bigl(\zeta+\norm{\nabla\Phi_T(q)}\bigr);
 \label{eq:scalar-chain-certificate}
\end{equation}
\item
\begin{equation}
 \Phi_T(0)-\inf_q\Phi_T(q)\le RT.
 \label{eq:scalar-chain-gap}
\end{equation}
In addition, if every $q_j\ge d$, then
$\nabla\Phi_T(q)=0$ and $\nabla^2\Phi_T(q)=0$.
\end{enumerate}
\end{lemma}

\begin{proof}
The first claim follows from boundedness of $s'$, $r$, and $r'$.  For the
second, let $j$ be the first index with $q_j\le b$.  Then
$s(q_{j-1})=1$ when $j>1$, and the same is true by convention for $j=1$.
Since $r'(q_j)=1$, \eqref{eq:scalar-chain-gradient} gives $g_j\ge1$.

We spell out the Hessian estimate.  Since $r$ is continuous and strictly
positive on $[-1,1]$, there is a universal $r_*>0$ such that
\begin{equation}
 r(u)\ge r_*,\qquad -1\le u\le1.
 \label{eq:reward-positive-lower}
\end{equation}
The only nonzero Hessian entries are
\begin{align*}
 \partial_{jj}^2\Phi_T
 &=-s(q_{j-1})r''(q_j)-s''(q_j)r(q_{j+1}),\\
 \partial_{j,j+1}^2\Phi_T
 &=-s'(q_j)r'(q_{j+1}).
\end{align*}
The reward estimate gives
\[
 s(q_{j-1})|r''(q_j)|
 \le C\Gamma_\zeta\bigl(\zeta+s(q_{j-1})r'(q_j)\bigr)
 \le C\Gamma_\zeta(\zeta+g_j).
\]
The flat-switch second-derivative estimate and
\eqref{eq:reward-positive-lower} give
\[
 |s''(q_j)|r(q_{j+1})
 \le C\Gamma_\zeta\bigl(\zeta+s'(q_j)\bigr)
 \le C\Gamma_\zeta(\zeta+g_j).
\]
For the off-diagonal entry, the left logarithmic-derivative bound for $s$
gives
\[
 s'(q_j)r'(q_{j+1})
 \le C\Gamma_\zeta\bigl(\zeta+s(q_j)r'(q_{j+1})\bigr)
 \le C\Gamma_\zeta(\zeta+g_{j+1}).
\]
Thus every absolute row sum of the tridiagonal Hessian is at most
$C\Gamma_\zeta(\zeta+\max_jg_j)$.  Since
$\max_jg_j\le\norm{\nabla\Phi_T(q)}$, the operator-norm bound
\eqref{eq:scalar-chain-certificate} follows.

Every summand in \eqref{eq:terminal-flat-scalar-chain} lies in $[-R,0]$,
which proves \eqref{eq:scalar-chain-gap}.  Finally, if every $q_j\ge d$, then
$r'(q_j)=r''(q_j)=s'(q_j)=s''(q_j)=0$, proving both gradient and Hessian
terminal flatness.
\end{proof}

\subsection{Block hiding and a public radial wall}

Let $D\ge2$, let $\theta_1,\ldots,\theta_T$ be unit vectors in mutually
separate copies of $\R^D$, and define
\begin{equation}
 \pi(z)=\frac{z}{\sqrt{1+\norm z^2}},
 \qquad
 q_j(z_j)=\inner{\theta_j}{\pi(z_j)}.
 \label{eq:block-compression}
\end{equation}
Then $q_j\in(-1,1)$.  Let
\begin{equation}
 \omega(r)=M\int_0^r\mathsf S(t-1)\,dt,
 \label{eq:radial-wall}
\end{equation}
where $M$ is a sufficiently large universal constant, and set
\begin{equation}
 \mathcal H_{\Theta,T}(z_1,\ldots,z_T)
 =\Phi_T(q_1(z_1),\ldots,q_T(z_T))
  +\sum_{j=1}^T\omega(\norm{z_j}).
 \label{eq:block-terminal-chain}
\end{equation}
The wall vanishes on the unit ball, has constant radial slope $M$ outside
the ball of radius two, and contains no hidden direction.

\begin{lemma}[Block-chain geometry and smoothness]
\label{lem:block-chain-geometry}
For a sufficiently large universal choice of $M$, there are universal
constants $c_g,C_g,C_H,C_\Delta>0$, independent of $T,D,\Theta$, such that:
\begin{enumerate}
\item if $q_j(z_j)\le b$ for at least one $j$, then
\begin{equation}
 \norm{\nabla\mathcal H_{\Theta,T}(z)}\ge c_g;
 \label{eq:block-gradient-obstruction}
\end{equation}
\item every hidden-chain block of $\nabla\mathcal H_{\Theta,T}$ has norm at
most $C_g$;
\item
\begin{equation}
 \norm{\nabla^2\mathcal H_{\Theta,T}(z)}_{\rm op}
 \le C_H\Gamma_\zeta
 \bigl(\zeta+\norm{\nabla\mathcal H_{\Theta,T}(z)}\bigr);
 \label{eq:block-generalized-certificate}
\end{equation}
\item
\begin{equation}
 \mathcal H_{\Theta,T}(0)-\inf_z\mathcal H_{\Theta,T}(z)
 \le C_\Delta T.
 \label{eq:block-chain-gap}
\end{equation}
\end{enumerate}
\end{lemma}

\begin{proof}
Fix one block and write $r=\norm z$, $e=z/r$ for $r>0$,
$c=\inner\theta e$, $u=\inner\theta z=rc$, and
$\rho=(1+r^2)^{-1/2}$.  Direct differentiation gives
\begin{align}
 \nabla q(z)
 &=\rho\theta-\rho^3uz,
 \label{eq:compression-gradient}\\
 \nabla^2q(z)
 &=-\rho^3(\theta z^\top+z\theta^\top+uI)
   +3\rho^5u zz^\top.
 \label{eq:compression-hessian}
\end{align}
Equivalently,
\begin{equation}
 \nabla q
 =\rho(\theta-ce)+\rho^3ce,
 \qquad
 \norm{\nabla q}^2
 =\rho^2(1-c^2)+\rho^6c^2.
 \label{eq:compression-radial-tangent}
\end{equation}
The singular values of $D\pi(z)$ are $\rho$ in tangential directions and
$\rho^3$ in the radial direction.  In particular,
\begin{equation}
 \norm{\nabla q}\ge5^{-3/2}\quad(r\le2),
 \qquad
 \norm{\nabla^2q}_{\rm op}\le6\rho^2.
 \label{eq:compression-quantitative-bounds}
\end{equation}
The Hessian estimate follows directly from
\eqref{eq:compression-hessian}, using $|u|\le r$,
$\rho r\le1$, and $\rho^3r^3\le1$.

Let $w(r)=\omega'(r)=M\mathsf S(r-1)$.  A block gradient has the form
\begin{equation}
 v=-g\nabla q+w(r)e,
 \qquad 0\le g\le G_0,
 \label{eq:block-gradient-two-parts}
\end{equation}
with the zero convention at $r=0$.  We claim that, for $M$ sufficiently
large,
\begin{equation}
 \norm v\ge c\bigl(g\norm{\nabla q}+w(r)\bigr)
 \label{eq:block-noncancellation}
\end{equation}
with a universal $c>0$.

If $r\le1$, then $w=0$ and the claim is immediate.  Suppose
$1\le r\le2$.  If $q\le0$, the radial component of $-\nabla q$ is outward,
so adding the outward wall cannot decrease either the chain norm or the wall
norm; hence \eqref{eq:block-noncancellation} follows.  If $q\ge d$, then
$r'(q)=s'(q)=0$ and consequently $g=0$.  In the remaining case
$0<q<d$, one has
\[
 |c|=\frac{|q|\sqrt{1+r^2}}r\le d\sqrt2<\frac12.
\]
Thus the tangential component of $-g\nabla q$ has norm at least
$c_0g\rho$, while $\norm{\nabla q}\le\rho$.  If $w\le2g\rho$, the
tangential component proves \eqref{eq:block-noncancellation}; if
$w>2g\rho$, the triangle inequality gives
$\norm v\ge w-g\norm{\nabla q}\ge w/2$.  Finally, when $r\ge2$ the wall
slope equals $M$.  If $q\ge d$, again $g=0$; otherwise choosing
$M\ge4G_0$ makes
$\norm v\ge M-G_0\norm{\nabla q}\ge3M/4$.  This proves the claim in all
cases.

The same cases, together with
\eqref{eq:compression-quantitative-bounds}, yield
\begin{equation}
 g\norm{\nabla^2q}_{\rm op}
 \le C\bigl(g\norm{\nabla q}+w(r)\bigr).
 \label{eq:compression-second-controlled}
\end{equation}
Indeed, for $r\le2$ the gradient is uniformly bounded below.  For $r\ge2$
and $0\le q<d$, the bound $|c|<1/2$ gives
$\norm{\nabla q}\ge c_0\rho$, while
$\norm{\nabla^2q}\le6\rho^2$.  For $r\ge2$ and $q<0$, the constant wall
slope dominates the bounded left side, and for $q\ge d$ one has $g=0$.

Apply \eqref{eq:block-noncancellation} to every block and put
$d_j=\norm{\nabla q_j}$ and $w_j=w(\norm{z_j})$.  Since the blocks are
orthogonal,
\begin{equation}
 \left(\sum_j(g_jd_j+w_j)^2\right)^{1/2}
 \le C\norm{\nabla\mathcal H_{\Theta,T}(z)}.
 \label{eq:coefficient-control-by-gradient}
\end{equation}
If some $q_j\le b$, Lemma~\ref{lem:scalar-chain-properties} gives a block
with $g_j\ge1$.  For that block either $r_j\le2$, in which case
$d_j\ge5^{-3/2}$, or $r_j>2$, in which case $w_j=M$.  The first assertion
therefore follows from \eqref{eq:block-noncancellation}.  The second follows
from $g_j\le G_0$ and $d_j\le1$.

For the Hessian, the lifted scalar-chain part is
\[
 J^\top\nabla^2\Phi_TJ
 +\operatorname{diag}(\partial_j\Phi_T\nabla^2q_j),
\]
where $J$ is block diagonal with blocks $\nabla q_j^\top$.  The proof of
Lemma~\ref{lem:scalar-chain-properties} gives
\begin{align*}
 |\partial_{jj}^2\Phi_T|
 &\le C\Gamma_\zeta(\zeta+g_j),\\
 |\partial_{j,j+1}^2\Phi_T|
 &\le C\Gamma_\zeta(\zeta+g_{j+1}).
\end{align*}
After multiplication by $d_jd_k$, every block-row sum is at most
\[
 C\Gamma_\zeta
 \left(\zeta+\max_\ell(g_\ell d_\ell+w_\ell)\right)
 \le C\Gamma_\zeta
 \left(\zeta+\norm{\nabla\mathcal H_{\Theta,T}(z)}\right),
\]
where we used \eqref{eq:coefficient-control-by-gradient}.  The second
composition term is controlled by
\eqref{eq:compression-second-controlled} in the same way.

For a radial wall, the Hessian eigenvalues are $w'(r)$ and $w(r)/r$.
They vanish for $r\le1$, while for $r\ge1$ the flat-switch estimate gives
\[
 |w'(r)|+\frac{w(r)}r
 \le C\Gamma_\zeta\bigl(\zeta+w(r)\bigr).
\]
The wall Hessian is block diagonal, so
\eqref{eq:coefficient-control-by-gradient} again gives
\eqref{eq:block-generalized-certificate}.

Finally, the walls are nonnegative and vanish at the origin, while the scalar
part lies in $[-RT,0]$.  Hence
\eqref{eq:block-chain-gap} holds with a universal $C_\Delta$.
\end{proof}

\subsection{Scaling into the target function class}

The next elementary lemma converts the pointwise Hessian certificate into
the exact finite-step condition of the problem.

\begin{lemma}[Vector differential certificate]
\label{lem:vector-differential-certificate}
Suppose a $C^2$ function $f$ satisfies
\begin{equation}
 \norm{\nabla^2f(x)}_{\rm op}
 \le\frac18\bigl(L_0+L_1\norm{\nabla f(x)}\bigr)
 \qquad\text{for every }x.
 \label{eq:vector-hessian-certificate}
\end{equation}
Then, whenever $\norm{x-y}\le1/L_1$,
\begin{equation}
 \norm{\nabla f(y)-\nabla f(x)}
 \le\bigl(L_0+L_1\norm{\nabla f(x)}\bigr)\norm{y-x}.
 \label{eq:vector-finite-step-certificate}
\end{equation}
\end{lemma}

\begin{proof}
Parametrize the segment by arclength and put
$Q(t)=\norm{\nabla f(x+tv)}+L_0/L_1$.  Its upper Dini derivative obeys
$Q'(t)\le(L_1/8)Q(t)$.  Hence
$Q(t)\le Q(0)e^{L_1t/8}$.  Integrating
\eqref{eq:vector-hessian-certificate} along a segment of length $d$ gives
\[
 \norm{\nabla f(y)-\nabla f(x)}
 \le\frac{L_0+L_1\norm{\nabla f(x)}}{L_1}
       (e^{L_1d/8}-1).
\]
For $0\le L_1d\le1$, $e^{L_1d/8}-1\le L_1d$, proving the claim.
\end{proof}

Assume for the moment that $L_1\eps\ge2L_0$ and define
\begin{equation}
 \zeta_\eps
 =\min\left\{e^{-4},\frac{c_gL_0}{16L_1\eps}\right\},
 \qquad
 \Gamma_\eps:=\Gamma_{\zeta_\eps}.
 \label{eq:coarse-chain-zeta}
\end{equation}
Choose
\begin{equation}
 \beta=\frac{L_1}{8C_H\Gamma_\eps},
 \qquad
 \alpha=\frac{8\eps}{c_g\beta},
 \label{eq:coarse-chain-scaling}
\end{equation}
and define
\begin{equation}
 f_{\Theta,T}(x)=\alpha\mathcal H_{\Theta,T}(\beta x).
 \label{eq:scaled-coarse-chain}
\end{equation}
By Lemma~\ref{lem:block-chain-geometry},
\begin{align*}
 \norm{\nabla^2 f_{\Theta,T}(x)}_{\rm op}
 &\le C_H\Gamma_\eps\alpha\beta^2\zeta_\eps
      +C_H\Gamma_\eps\beta\norm{\nabla f_{\Theta,T}(x)}\\
 &\le\frac18L_0+\frac18L_1\norm{\nabla f_{\Theta,T}(x)}.
\end{align*}
Thus Lemma~\ref{lem:vector-differential-certificate} proves the required
relaxed smoothness.  Moreover,
\begin{equation}
 \text{if some }q_j\le b,
 \qquad
 \norm{\nabla f_{\Theta,T}(x)}\ge8\eps.
 \label{eq:scaled-gradient-obstruction}
\end{equation}
The scaled gap is at most $\alpha C_\Delta T$.  Consequently, for a
sufficiently small universal $c_T>0$, the choice
\begin{equation}
 T=\left\lfloor
 c_T\frac{B_\eps}{\Gamma_\eps}
 \right\rfloor
 \label{eq:coarse-chain-length}
\end{equation}
ensures
$f_{\Theta,T}(0)-\inf f_{\Theta,T}\le\Delta$.

\subsection{A fixed fresh rare-reveal oracle}

The oracle must be defined pointwise for every query, including queries that
deactivate a previously active block.  For this reason we use the maximum
active index rather than the length of a currently active contiguous prefix.
For a base point $z=(z_1,\ldots,z_T)$ define
\begin{equation}
 k_\Theta(z)
 =\max\bigl(\{0\}\cup\{j:q_j(z_j)>a\}\bigr).
 \label{eq:max-active-index}
\end{equation}

\begin{lemma}[One-new-direction decomposition]
\label{lem:one-new-direction-decomposition}
For every $z$, if $k=k_\Theta(z)<T$, then
\begin{equation}
 \nabla\mathcal H_{\Theta,T}(z)
 =D_{\Theta_{\le k}}(z)+h_{\Theta,k+1}(z),
 \label{eq:one-new-direction-base}
\end{equation}
where $D_{\Theta_{\le k}}$ depends only on
$\theta_1,\ldots,\theta_k$, no direction with index at least $k+2$ appears,
and
\begin{equation}
 \sup_{T,D,\Theta,z,k}\norm{h_{\Theta,k+1}(z)}\le C_h
 \label{eq:rare-vector-bounded}
\end{equation}
for a universal constant $C_h$.  If $k=T$, take
$h_{\Theta,T+1}=0$ and $D_{\Theta_{\le T}}=\nabla\mathcal H_{\Theta,T}$.
\end{lemma}

\begin{proof}
The hidden-chain contribution in block $j$ is
$-g_j(q)\nabla q_j$, where
\[
 g_j(q)=s(q_{j-1})r'(q_j)+s'(q_j)r(q_{j+1}).
\]
If $j\ge k+2$, then $q_{j-1}\le a$ and $q_j\le a$.  Hence
$s(q_{j-1})=s'(q_j)=0$, so $g_j=0$.  Thus no direction beyond $k+1$ appears.

For $1\le k<T$, split the only two blocks that may depend on
$\theta_{k+1}$ as
\begin{align}
 [h_{\Theta,k+1}(z)]_k
 &=-s'(q_k)\bigl(r(q_{k+1})-r(0)\bigr)\nabla q_k,
 \label{eq:rare-vector-prev-block}\\
 [h_{\Theta,k+1}(z)]_{k+1}
 &=-s(q_k)r'(q_{k+1})\nabla q_{k+1}.
 \label{eq:rare-vector-new-block}
\end{align}
All remaining blocks of $h_{\Theta,k+1}$ are zero.  The baseline term
$-s'(q_k)r(0)\nabla q_k$, every block of index at most $k-1$, and the public
radial wall are placed in $D_{\Theta_{\le k}}$.  Because
$q_{k+1}\le a<b$, one has $s'(q_{k+1})=0$ and $r'(q_{k+1})=1$, so this split
is exact.  For $k=0$, all coordinates lie on the lower gate plateau and the
only hidden-chain block is $-\nabla q_1$; take this as $h_{\Theta,1}$.
Uniform boundedness of $s'$, $r$, $r'$, and $\nabla q_j$ proves
\eqref{eq:rare-vector-bounded}.
\end{proof}

In scaled coordinates, with $z=\beta x$, put
\begin{equation}
 H_{\Theta,k+1}(x)=\alpha\beta h_{\Theta,k+1}(\beta x),
 \qquad
 \norm{H_{\Theta,k+1}(x)}\le C_*\eps,
 \label{eq:scaled-rare-vector}
\end{equation}
where $C_*=8C_h/c_g$.  Define
\begin{equation}
 p=
 \begin{cases}
  1,&\sigma=0,\\
  \min\left\{1,\dfrac{C_*^2\eps^2}{\sigma^2}\right\},&\sigma>0.
 \end{cases}
 \label{eq:rare-probability}
\end{equation}
For a fresh $\xi\sim\operatorname{Bernoulli}(p)$ and
$k=k_\Theta(\beta x)$, define the fixed oracle
\begin{equation}
 G_\Theta(x;\xi)
 =\alpha\beta D_{\Theta_{\le k}}(\beta x)
  +\frac{\xi}{p}H_{\Theta,k+1}(x).
 \label{eq:rare-reveal-oracle}
\end{equation}
When $k=T$, the second term is zero and the oracle returns the exact gradient.
Lemma~\ref{lem:one-new-direction-decomposition} gives, pointwise for every
query $x$,
\begin{equation}
 \E_\xi[G_\Theta(x;\xi)]=\nabla f_{\Theta,T}(x),
 \qquad
 \E_\xi\norm{G_\Theta(x;\xi)-\nabla f_{\Theta,T}(x)}^2
 =\frac{1-p}{p}\norm{H_{\Theta,k+1}(x)}^2
 \le\sigma^2.
 \label{eq:rare-oracle-validity}
\end{equation}
Moreover,
\begin{equation}
 p^{-1}=\Theta(P_\eps)
 \label{eq:p-is-P}
\end{equation}
up to universal constants.  Notice that this is a single fixed oracle kernel:
$k_\Theta(\beta x)$ is a deterministic function of the instance and query,
and every call uses a fresh independent Bernoulli variable.

\subsection{No parallel bypass and the product lower bound}

Draw $\theta_1,\ldots,\theta_T$ independently and uniformly from the unit
sphere in $\R^D$.  For the analysis, strengthen a successful oracle reply by
also revealing the corresponding new direction in full.  Let $\ell_0=0$ and
let $\ell_t$ be the largest initial segment of directions revealed after the
first $t$ calls.  If at call $t$ one has
$k_\Theta(\beta x_t)=\ell_{t-1}<T$ and $\xi_t=1$, set
$\ell_t=\ell_{t-1}+1$; otherwise set $\ell_t=\ell_{t-1}$.  This extra reveal
can only help the algorithm.

\begin{lemma}[Blockwise no-bypass and monotone exposure]
\label{lem:block-no-bypass}
Fix a possibly randomized algorithm making at most $Q$ oracle calls and then
returning one output.  If
\begin{equation}
 D\ge C a^{-2}\log\bigl(16T(Q+1)\bigr),
 \label{eq:block-dimension-choice}
\end{equation}
then, with probability at least $7/8$, the following event holds
simultaneously at every query and at the output:
\begin{equation}
 q_j(\beta x_{t,j})\le a
 \quad\text{for every }j>\ell_{t-1},
 \qquad
 q_j(\beta\widehat x_j)\le a
 \quad\text{for every }j>\ell_Q.
 \label{eq:no-bypass-event}
\end{equation}
On this event, one Bernoulli success increases $\ell_t$ by at most one, and a
non-success reveals no new direction.
\end{lemma}

\begin{proof}
We use a truncated interaction so that no independence assertion is made
after conditioning on the absence of an earlier bypass.  Fix the algorithm's
private seed, and couple all interactions using the same independent
Bernoulli variables and hidden directions.  We may additionally disclose
each Bernoulli variable after its call and disclose a new direction whenever
the augmented revelation rule calls for it.  These disclosures only
strengthen the algorithm.

In the truncated interaction, maintain a revealed prefix of length
$\widetilde\ell$, initially zero.  At a query $z=\beta x$, form
\[
 \widetilde k
 =\max\bigl(\{0\}\cup
       \{j\le\widetilde\ell:q_j(z_j)>a\}\bigr).
\]
This index uses revealed directions only.  The expressions for
$D_{\Theta_{\le k}}$ and $h_{\Theta,k+1}$ in the proof of
Lemma~\ref{lem:one-new-direction-decomposition} also define algebraic
extensions outside the region $k=k_\Theta(z)$: retain the public wall,
the displayed prefix terms and the displayed successor correction, and
omit every further hidden-chain contribution.  Use these extensions with
$k=\widetilde k$ to form the reply.  If
$\widetilde k<\widetilde\ell$, all required directions are already known.
If $\widetilde k=\widetilde\ell<T$, a Bernoulli failure returns only
$\alpha\beta D_{\Theta_{\le\widetilde\ell}}(z)$; a success reads and
discloses $\theta_{\widetilde\ell+1}$, returns the corresponding successor
correction, and increases the prefix length by one.  When the prefix has
length $T$, all directions are known and no further disclosure is needed.

By induction, this simulated history reads only its revealed prefix.
In particular, the decision to read a successor direction is made before
that direction is read.  Conditional on the simulated history before a
call, all directions outside its prefix remain independent uniform sphere
vectors.  This statement concerns the unconditioned truncated interaction;
it does not condition on any alignment event.  For each as-yet unread
direction, the standard spherical-cap bound therefore gives
\[
 \Pr\!\left(
   \inner{\theta_j}{\pi(\beta\widetilde x_{t,j})}>a
   \,\middle|\,\widetilde{\mathcal F}_{t-1}
 \right)
 \le e^{-cDa^2}
 \qquad(j>\widetilde\ell_{t-1}),
\]
after adjusting the universal constant for $D\ge2$.
The same bound holds for the simulated final output and each direction
outside its final prefix.  A union bound over at most $T(Q+1)$ pairs shows
that the probability of any such simulated alignment exceeding $a$ is at
most $1/8$ under~\eqref{eq:block-dimension-choice}.

On the complementary event the truncated and true augmented interactions
coincide, pathwise.  Indeed, assume their histories agree before a query.
All directions outside the common revealed prefix then have alignment at
most $a$, so the true maximum active index equals $\widetilde k$.
Lemma~\ref{lem:one-new-direction-decomposition} makes their oracle replies
identical, including the successor correction when it is sampled.  Their
disclosures, next queries, and prefix lengths consequently agree as well.
Induction gives equality through all calls and at the final output.  This
proves~\eqref{eq:no-bypass-event}, together with the asserted one-direction
exposure rule.  Averaging over the private seed completes the proof.
\end{proof}

\subsection{Proof of Theorem~\ref{thm:stochastic-coarse-product}}

\begin{proof}
Choose the universal constant in $B_\eps\ge C\Xi_\eps$ large enough
that $T\ge32$. Consider first an algorithm with a deterministic budget
\begin{equation}
 Q=\left\lfloor\frac{T}{16p}\right\rfloor.
 \label{eq:coarse-fixed-budget-Q}
\end{equation}
Choose $D$ according to Lemma~\ref{lem:block-no-bypass}.  The number of
Bernoulli successes among the $Q$ calls has expectation at most $T/16$, so
Markov's inequality gives probability at most $1/16$ that there are $T$ or
more successes.  Lemma~\ref{lem:block-no-bypass} fails with probability at
most $1/8$.  Hence, with probability at least $13/16$, the output has at
least one unrevealed block with $q_j\le a<b$.  By
\eqref{eq:scaled-gradient-obstruction},
\begin{equation}
 \norm{\nabla f_{\Theta,T}(\widehat x)}\ge8\eps
 \label{eq:coarse-product-failure-gradient}
\end{equation}
on that event.  Averaging over the independent random directions and then
fixing one realization shows that every randomized fixed-budget algorithm
has a fixed admissible instance on which its expected output gradient is
larger than $\eps$.

Now let an algorithm have worst-case expected budget $q$.  Cap it after the
$Q$ calls in \eqref{eq:coarse-fixed-budget-Q}.  The fixed-budget argument
selects one realization $\Theta_\star$ for which the capped algorithm obeys
\eqref{eq:coarse-product-failure-gradient} with probability at least
$13/16$.  If $q\le Q/32$, then on that same instance
$\Pr(S>Q)\le1/32$.  The original and capped outputs agree on
$\{S\le Q\}$, and therefore the original output has gradient at least
$8\eps$ with probability at least $25/32$, contradicting the target
criterion.  Thus
\[
 q\ge c\frac{T}{p}
 =\Omega\left(\frac{B_\eps P_\eps}{\Gamma_\eps}\right).
\]
Finally, \eqref{eq:coarse-chain-zeta} implies
$\Gamma_\eps=O(\Xi_\eps)$ under
$L_1\eps\ge2L_0$.  Giving every node the same objective and independent
copies of \eqref{eq:rare-reveal-oracle} converts the centralized query lower
bound into the total decentralized local-call lower bound without changing
any parameter.
\end{proof}

\section{A one-dimensional deterministic coarse lower bound}
\label{sec:coarse-reset}

The main lower bounds are allowed to increase the ambient dimension. This
appendix records a complementary one-dimensional obstruction to coarse
optimization. It uses identical local objectives and exact gradients, and
establishes the $B_\eps$ scale up to a logarithm without a high-dimensional
hiding argument. The oracle is gradient-only, as in the main model.

The basic device is a periodic derivative which is equal to a fixed positive
slope in every cell except one hidden cell.  Inside that cell it descends
exponentially, spends enough time on a negative plateau to reset the total
integral, and then returns exponentially to the common positive slope.  The
exponential ramps are chosen so that their derivative is proportional to
$L_0+L_1|f'|$; consequently the construction obeys the stated finite-step
relaxed-smoothness inequality exactly.

\begin{lemma}[Differential certificate for relaxed smoothness]
\label{lem:coarse-differential-certificate}
Let $p:\R\to\R$ be locally absolutely continuous and suppose that, almost
everywhere,
\begin{equation}
 |p'(x)|\le \frac14\bigl(L_0+L_1|p(x)|\bigr).
 \label{eq:coarse-differential-bound}
\end{equation}
If $f'(x)=p(x)$, then for all $x,y$ with $|x-y|\le1/L_1$,
\begin{equation}
 |f'(y)-f'(x)|
 \le \bigl(L_0+L_1|f'(x)|\bigr)|y-x|.
 \label{eq:coarse-exact-relaxed}
\end{equation}
\end{lemma}

\begin{proof}
Put $c=L_0/L_1$ and $q=|p|+c$.  Along either orientation of the segment from
$x$ to $y$, \eqref{eq:coarse-differential-bound} and Gronwall's inequality
give
\[
 q(x+s)\le q(x)e^{L_1s/4},
 \qquad 0\le s\le |y-x|.
\]
Therefore, writing $d=|y-x|$,
\[
 |p(y)-p(x)|
 \le \frac{L_1}{4}\int_0^d q(x+s)\,ds
 \le q(x)\bigl(e^{L_1d/4}-1\bigr).
\]
For $0\le u\le1$, one has $e^{u/4}-1\le u$.  Since $L_1d\le1$, the last
display is at most
$L_1q(x)d=(L_0+L_1|p(x)|)d$.
\end{proof}

Define the logarithmic loss
\begin{equation}
 \Lambda_\eps
 :=1+\log\left(1+4\frac{L_1\eps}{L_0}\right)
      +\log\left(1+8(1+B_\eps)\frac{L_1\eps}{L_0}\right).
 \label{eq:coarse-log-loss}
\end{equation}

\begin{theorem}[Hidden reset-cycle lower bound]
\label{thm:coarse-reset-lower}
There are universal constants $c,C>0$ with the following property.  Suppose
\begin{equation}
 L_1\eps\ge L_0,
 \qquad
 B_\eps\ge C\Lambda_\eps.
 \label{eq:coarse-reset-regime}
\end{equation}
Then every randomized adaptive decentralized algorithm has an admissible
one-dimensional exact-gradient instance for which
\begin{equation}
 \boxed{
 \E S\ge c\frac{B_\eps}{\Lambda_\eps}.}
 \label{eq:coarse-B-lower}
\end{equation}
All local objectives in the instance are identical.  Hence the construction
has zero heterogeneity and its deterministic oracle is an admissible subclass
for every stated value of $H$ and $\sigma$.
\end{theorem}

\begin{proof}
Fix an integer $K\ge16$, to be chosen below, and set
\begin{equation}
 \alpha=4\eps,
 \qquad
 \kappa=\frac14,
 \qquad
 c_0=\frac{L_0}{L_1},
 \qquad
 M=2K\alpha.
 \label{eq:coarse-pulse-scales}
\end{equation}
The assumption $L_1\eps\ge L_0$ implies
$r:=\alpha/c_0\ge4$.  Define
\begin{align}
 \ell_+
 &=\frac1{\kappa L_1}\log\left(1+\frac{\alpha}{c_0}\right),
 &
 \ell_-
 &=\frac1{\kappa L_1}\log\left(1+\frac{M}{c_0}\right),
 &
 \ell&=\ell_++\ell_- .
 \label{eq:coarse-ramp-lengths}
\end{align}
On $[0,\ell]$ introduce the descending profile
\begin{equation}
 d(t)=
 \begin{cases}
  (\alpha+c_0)e^{-\kappa L_1t}-c_0,
       &0\le t\le\ell_+,\\[1mm]
  c_0-c_0e^{\kappa L_1(t-\ell_+)},
       &\ell_+\le t\le\ell.
 \end{cases}
 \label{eq:coarse-descending-profile}
\end{equation}
It satisfies $d(0)=\alpha$, $d(\ell_+)=0$, $d(\ell)=-M$, and
\begin{equation}
 |d'(t)|=\kappa\bigl(L_0+L_1|d(t)|\bigr)
 \quad\text{almost everywhere.}
 \label{eq:coarse-ramp-ode}
\end{equation}
Its signed integral is
\begin{equation}
 D:=\int_0^\ell d(t)\,dt
 =\frac1{\kappa L_1}
 \left[
  \alpha-M+c_0\log\frac{c_0+M}{c_0+\alpha}
 \right].
 \label{eq:coarse-ramp-integral}
\end{equation}
Choose the negative-plateau length
\begin{equation}
 h=
 \frac{2\alpha(K-1)\ell+2D}
      {M-\alpha(K-1)}.
 \label{eq:coarse-plateau-length}
\end{equation}
The denominator is $\alpha(K+1)$.  Moreover,
\begin{equation}
 -\frac{M}{\kappa L_1}\le D\le\frac{\alpha}{\kappa L_1}.
 \label{eq:coarse-D-bounds}
\end{equation}
Indeed, the logarithm in \eqref{eq:coarse-ramp-integral} is nonnegative,
while
$c_0\log((c_0+M)/(c_0+\alpha))\le
c_0\log(1+M/c_0)\le M$.
Since $M/c_0=2Kr\ge128$, \eqref{eq:coarse-D-bounds} gives
\[
 2\alpha(K-1)\ell+2D
 \ge
 \frac{2\alpha}{\kappa L_1}
 \left[(K-1)\log(1+2Kr)-2K\right]
 \ge0.
\]
Thus $h\ge0$.  The upper bound in \eqref{eq:coarse-D-bounds} also gives
\begin{equation}
 h\le2\ell+\frac2{\kappa L_1},
 \qquad
 w:=2\ell+h
 \le4\ell+\frac2{\kappa L_1}.
 \label{eq:coarse-cell-width}
\end{equation}

Define a continuous pulse $\pi:[0,w]\to[-M,\alpha]$ by
\begin{equation}
 \pi(t)=
 \begin{cases}
  d(t),&0\le t\le\ell,\\
  -M,&\ell\le t\le\ell+h,\\
  d(w-t),&\ell+h\le t\le w.
 \end{cases}
 \label{eq:coarse-pulse}
\end{equation}
The pulse begins and ends at $\alpha$, and its integral is $2D-Mh$.
Let ${\mathsf P}=Kw$ and divide one period into cells
$I_j=[jw,(j+1)w)$, $j=0,\ldots,K-1$.  For a hidden index
$J\in\{0,\ldots,K-1\}$, define the ${\mathsf P}$-periodic derivative
$p_J$ to equal the translated pulse $\pi$ on $I_J$ and the constant
$\alpha$ on every other cell.  By \eqref{eq:coarse-plateau-length},
\begin{equation}
 \int_0^{\mathsf P}p_J(t)\,dt
 =\alpha(K-1)w+2D-Mh=0.
 \label{eq:coarse-zero-mean}
\end{equation}
Consequently
\begin{equation}
 f_J(x):=\int_0^x p_J(t)\,dt
 \label{eq:coarse-objective}
\end{equation}
is continuously differentiable, periodic, and bounded below.  The positive
part of $p_J$ is everywhere at most $\alpha$.  Since the signed integral over
a period is zero, the oscillation of $f_J$ is at most the total positive
variation, hence
\begin{equation}
 f_J(0)-\inf f_J
 \le \operatorname{osc}(f_J)
 \le \alpha {\mathsf P}=\alpha Kw.
 \label{eq:coarse-gap-bound}
\end{equation}
The derivative $p_J$ is locally absolutely continuous and satisfies
\eqref{eq:coarse-differential-bound}: equality up to the factor $\kappa$
holds on the two exponential ramps, and the derivative vanishes on all
plateaux.  Lemma~\ref{lem:coarse-differential-certificate} therefore proves
the required relaxed smoothness globally.

We now choose $K$ while respecting the gap budget.  Let
\begin{equation}
 K=\left\lfloor c_K\frac{B_\eps}{\Lambda_\eps}\right\rfloor
 \label{eq:coarse-K-choice}
\end{equation}
for a sufficiently small universal $c_K>0$.  Assumption
\eqref{eq:coarse-reset-regime}, after increasing $C$, ensures $K\ge16$.
Also $K\le B_\eps$, and therefore
$M/c_0=8K L_1\eps/L_0\le8B_\eps L_1\eps/L_0$.
Equations \eqref{eq:coarse-ramp-lengths} and
\eqref{eq:coarse-cell-width} yield
\begin{equation}
 w\le \frac{C_w\Lambda_\eps}{L_1}
 \label{eq:coarse-width-log}
\end{equation}
for a universal $C_w$.  Hence
\[
 \alpha Kw
 \le 4C_wc_K\frac{\eps B_\eps}{L_1}
 =4C_wc_K\Delta.
\]
Taking $c_K\le(4C_w)^{-1}$ makes
\eqref{eq:coarse-gap-bound} at most $\Delta$.

It remains to prove the query lower bound.  Give every agent the same local
objective $f_J$ and the exact-gradient oracle $p_J$.  Consider first a
deterministic protocol using at most $n$ total local queries, and expose the
hypothetical transcript in which every oracle response equals $\alpha$.
The periodic extension identifies every point with one residue cell modulo
$\mathsf P$.  A query therefore tests at most one of the $K$ residue classes;
boundary points return $\alpha$ and may be assigned arbitrarily.  Thus the
query sequence touches at most $n$ candidate residue classes.  Fix one agent
and include the residue class containing its output, for a total of at most
$n+1$ classes.  For every hidden index $J$ outside this set, the actual
transcript under $f_J$ is exactly the hypothetical all-$\alpha$ transcript,
by induction over the adaptive queries, and the output lies outside every
translate of the hidden cell $I_J$.
Therefore
\begin{equation}
 |f'_J(\widehat x)|=\alpha=4\eps.
 \label{eq:coarse-unqueried-output}
\end{equation}
Under a uniform hidden index, every deterministic $n$-query protocol thus
satisfies
\begin{equation}
 \Pr_J\bigl(|f'_J(\widehat x)|=4\eps\bigr)
 \ge1-\frac{n+1}{K}.
 \label{eq:coarse-fixed-budget-failure}
\end{equation}
The same statement holds for randomized fixed-budget protocols after
conditioning on their private seed and averaging.

Finally let a randomized protocol have worst-case expected total call count
$q$.  Cap it after $n=\lfloor K/4\rfloor$ calls, using any
transcript-measurable fallback output.  By
\eqref{eq:coarse-fixed-budget-failure}, averaging over the hidden index and
then fixing one index $J_\star$ gives failure probability at least $11/16$
for the capped protocol.  If $q\le K/128$, Markov's inequality on that same
instance gives
\[
 \Pr_{J_\star}(S>n)
 \le\frac{q}{n+1}\le\frac1{32}.
\]
The capped and original outputs agree on $\{S\le n\}$, so the original
protocol has
\[
 \Pr_{J_\star}\bigl(|f'_{J_\star}(\widehat x)|=4\eps\bigr)
 \ge\frac{11}{16}-\frac1{32}
 =\frac{21}{32}.
\]
Its expected gradient norm is therefore larger than $\eps$, a contradiction.
Thus $q>K/128=\Omega(B_\eps/\Lambda_\eps)$.
\end{proof}

\end{document}
